\ifdefined\BookMaster
\def\BookEnd{}
\else
\documentclass[a4paper,12pt,fleqn,openany,twoside,fontset=windows]{ctexbook}

\usepackage[fleqn]{amsmath}
\usepackage{amssymb,amsthm,mathrsfs}
\usepackage{geometry}
\usepackage[dvipsnames]{xcolor}
\usepackage{graphicx}
\usepackage{tikz}
\usetikzlibrary{calc}
\usepackage[most]{tcolorbox}
\usepackage{fancyhdr}
\usepackage{titlesec}
\usepackage{tocloft}
\usepackage{enumitem}
\usepackage{microtype}
\usepackage{pifont}
\usepackage{hyperref}
\usepackage{romannum}
\usepackage{mathtools}

\definecolor{bookblue}{HTML}{264653}
\definecolor{bookteal}{HTML}{4F7C82}
\definecolor{booklight}{HTML}{EDF4F3}
\definecolor{bookgray}{HTML}{5E686B}

\geometry{
    inner=2.7cm,
    outer=2.5cm,
    top=2.7cm,
    bottom=2.7cm,
    headheight=18pt,
    headsep=18pt,
    footskip=25pt
}
\setlength{\mathindent}{0pt}
\setlength{\parindent}{5pt}
\setlength{\parskip}{3pt plus 1pt minus 1pt}
\linespread{1.25}
\allowdisplaybreaks[3]
\AtBeginDocument{%
    \setlength{\parindent}{5pt}
    \setlength{\abovedisplayskip}{8pt plus 2pt minus 2pt}
    \setlength{\belowdisplayskip}{8pt plus 2pt minus 2pt}
    \setlength{\abovedisplayshortskip}{5pt plus 1pt minus 1pt}
    \setlength{\belowdisplayshortskip}{5pt plus 1pt minus 1pt}
}

\hypersetup{
    unicode=true,
    colorlinks=true,
    linkcolor=bookblue,
    citecolor=bookblue,
    urlcolor=bookteal,
    bookmarksopen=true,
    pdfauthor={Chen},
    pdftitle={数学分析习题整理}
}

\renewcommand{\chaptermark}[1]{%
    \edef\@tmp{第\thechapter 章\quad #1}%
    \expandafter\markboth\expandafter{\@tmp}{}%
}
\fancypagestyle{main}{%
    \fancyhf{}
    \fancyhead[LE]{\small\color{bookgray}\nouppercase{\leftmark}}
    \fancyhead[RO]{\small\color{bookgray}数学分析习题整理}
    \fancyfoot[C]{\color{bookblue}\thepage}
    \fancyfoot[R]{\small\bfseries Chen\;\; github.com/mathlover-1/math-analysis}
    \renewcommand{\headrulewidth}{0.4pt}
    \renewcommand{\footrulewidth}{0pt}
}
\fancypagestyle{plain}{%
    \fancyhf{}
    \fancyfoot[C]{\color{bookblue}\thepage}
    \fancyfoot[R]{\small\bfseries Chen\;\; github.com/mathlover-1/math-analysis}
    \renewcommand{\headrulewidth}{0pt}
    \renewcommand{\footrulewidth}{0pt}
}
\pagestyle{main}

\titleformat{\chapter}[display]
  {\normalfont\sffamily\bfseries\color{bookblue}}
  {\raggedleft\fontsize{46}{50}\selectfont\thechapter}
  {-18pt}
  {\titlerule[1.2pt]\vspace{10pt}\filright\Huge}
\titlespacing*{\chapter}{0pt}{-15pt}{36pt}

\renewcommand{\contentsname}{目\quad 录}
\renewcommand{\cfttoctitlefont}{\hfill\sffamily\bfseries\Huge\color{bookblue}}
\renewcommand{\cftaftertoctitle}{\hfill}
\setlength{\cftbeforetoctitleskip}{5pt}
\setlength{\cftaftertoctitleskip}{28pt}
\setlength{\cftbeforechapskip}{12pt}
\renewcommand{\cftchapfont}{\sffamily\bfseries\large\color{bookblue}}
\renewcommand{\cftchappagefont}{\sffamily\bfseries\large\color{bookblue}}
\renewcommand{\cftchapleader}{\color{bookteal}\cftdotfill{2.5}}

\newtcolorbox{ti}{
    enhanced,
    breakable,
    colback=booklight!55!white,
    colframe=bookteal!35!white,
    boxrule=0.45pt,
    boxsep=0pt,
    left=10pt,
    right=10pt,
    top=9pt,
    bottom=9pt,
    before skip=14pt,
    after skip=14pt,
    arc=2pt,
    outer arc=2pt,
    before upper={\parindent=0pt}
}

\newenvironment{booknotebody}{%
    \begin{center}
    \begin{minipage}{0.84\textwidth}
    \songti\zihao{-4}\linespread{1.45}\selectfont
    \setlength{\parindent}{2\ccwd}
    \setlength{\parskip}{0.75em}
}{%
    \end{minipage}
    \end{center}
}

\newenvironment{booknotelongbody}{%
    \begingroup
    \songti\zihao{-4}\linespread{1.45}\selectfont
    \setlength{\parindent}{2\ccwd}
    \setlength{\parskip}{0.75em}
    \begin{list}{}{%
        \setlength{\leftmargin}{0.08\textwidth}
        \setlength{\rightmargin}{0.08\textwidth}
        \setlength{\topsep}{0pt}
        \setlength{\partopsep}{0pt}
        \setlength{\parsep}{0pt}
        \setlength{\itemsep}{0pt}
    }
    \item\relax
}{%
    \end{list}
    \endgroup
}

\fancypagestyle{plainnowm}{%
    \fancyhf{}
    \fancyfoot[C]{\color{bookblue}\thepage}
    \renewcommand{\headrulewidth}{0pt}
    \renewcommand{\footrulewidth}{0pt}
}

\newcommand{\booknotetitle}[2][plain]{%
    \clearpage
    \phantomsection
    \addcontentsline{toc}{chapter}{#2}
    \markboth{#2}{}
    \thispagestyle{#1}
    \vspace*{1.2cm}
    \begin{center}
        {\sffamily\heiti\bfseries\fontsize{26}{34}\selectfont\color{bookblue}#2\par}
        \vspace{0.8em}
        {\color{bookteal}\rule{4.5cm}{0.8pt}\par}
    \end{center}
    \vspace{0.6cm}
}

\renewcommand{\proofname}{证明}
\renewcommand{\qedsymbol}{\color{bookteal}\ensuremath{\blacksquare}}

\title{数学分析习题整理}
\author{Chen}
\date{}

\makeatletter
% ==================== 旧版封面/封底设计备份 ====================
% 2026-08 重设计,旧代码用 \iffalse...\fi 封存,新版见下方。
\iffalse
\renewcommand{\maketitle}{%
    \begin{titlepage}
        \thispagestyle{empty}
        \setcounter{page}{0}
        \noindent\hspace*{-\parindent}\begin{tikzpicture}[remember picture,overlay]
            \fill[bookblue] (current page.north west) rectangle ([yshift=-4.2cm]current page.north east);
            \fill[booklight] ([yshift=-4.2cm]current page.north west) rectangle (current page.south east);
            \draw[bookteal!55,line width=0.7pt]
                ([xshift=2.4cm,yshift=4.3cm]current page.south west)
                .. controls +(2.5cm,1.8cm) and +(-2.8cm,-1.3cm) ..
                ([xshift=-2.4cm,yshift=7.2cm]current page.south east);
            \draw[bookteal!35,line width=0.7pt]
                ([xshift=2.5cm,yshift=6.2cm]current page.south west)
                .. controls +(3.0cm,-1.2cm) and +(-3.0cm,1.6cm) ..
                ([xshift=-2.3cm,yshift=3.8cm]current page.south east);
        \end{tikzpicture}
        \vspace*{5.6cm}
        \noindent\color{bookblue}\rule{\textwidth}{1.4pt}\par
        \vspace{8pt}
        {\centering\sffamily\bfseries\fontsize{32}{42}\selectfont\color{bookblue}\@title\par}
        \vspace{8pt}
        \noindent\color{bookteal}\rule{\textwidth}{0.6pt}\par
        \vspace{2.6cm}
        {\centering\Large\color{bookgray}\@author\par}
        \vfill
        {\centering\color{bookteal}\large
            $\displaystyle \lim_{n\to\infty} a_n$\qquad
            $\displaystyle \int_a^b f(x)\,\mathrm{d}x$\qquad
            $\displaystyle \nabla f(x,y)$\par
        }
        \vspace*{1.5cm}
    \end{titlepage}
}

\newcommand{\makebackcover}{%
    \begingroup
    \let\cleardoublepage\clearpage
    \begin{titlepage}
        \thispagestyle{empty}
        \noindent\hspace*{-\parindent}\begin{tikzpicture}[remember picture,overlay]
            \fill[bookblue] (current page.north west) rectangle (current page.south east);
            \fill[booklight,rounded corners=7pt]
                ([xshift=2.2cm,yshift=-3.0cm]current page.north west)
                rectangle
                ([xshift=-2.2cm,yshift=3.0cm]current page.south east);
            \fill[bookteal!80]
                (current page.north west) --
                ([xshift=5.3cm]current page.north west) --
                ([xshift=2.1cm,yshift=-4.4cm]current page.north west) --
                ([yshift=-6.2cm]current page.north west) -- cycle;
            \fill[bookteal!55]
                (current page.south east) --
                ([xshift=-5.8cm]current page.south east) --
                ([xshift=-2.4cm,yshift=4.2cm]current page.south east) --
                ([yshift=6.0cm]current page.south east) -- cycle;
            \foreach \radius in {0.75,1.15,1.55,1.95} {
                \draw[bookteal!55,line width=0.55pt]
                    ([xshift=-2.9cm,yshift=-4.4cm]current page.north east)
                    circle[radius=\radius cm];
            }
            \foreach \radius in {0.65,1.05,1.45} {
                \draw[booklight!65,line width=0.55pt]
                    ([xshift=2.4cm,yshift=3.6cm]current page.south west)
                    circle[radius=\radius cm];
            }
            \fill[bookteal!16,rounded corners=3pt]
                ([xshift=3.5cm,yshift=-7.1cm]current page.north west)
                rectangle
                ([xshift=-5.0cm,yshift=7.2cm]current page.south east);
            \fill[bookteal!30,rounded corners=3pt]
                ([xshift=4.4cm,yshift=-8.0cm]current page.north west)
                rectangle
                ([xshift=-4.1cm,yshift=8.1cm]current page.south east);
            \begin{scope}[shift={([xshift=-0.75cm]current page.center)}]
                \clip[rounded corners=3pt] (-5.25,-4.75) rectangle (5.25,4.75);
                \foreach \v in {-3.6,-3.0,...,3.6} {
                    \draw[bookteal!48,line width=0.45pt,opacity=0.72]
                        plot[domain=-3.6:3.6,samples=45,smooth,variable=\u]
                        ({\u+0.55*\v},{0.28*\u-0.35*\v+0.11*(\u*\u-\v*\v)});
                }
                \foreach \u in {-3.6,-3.0,...,3.6} {
                    \draw[bookblue!42,line width=0.42pt,opacity=0.62]
                        plot[domain=-3.6:3.6,samples=45,smooth,variable=\v]
                        ({\u+0.55*\v},{0.28*\u-0.35*\v+0.11*(\u*\u-\v*\v)});
                }
                \draw[bookteal!70,line width=0.75pt,opacity=0.82]
                    plot[domain=-3.6:3.6,samples=55,smooth,variable=\u]
                    ({\u},{0.28*\u+0.11*\u*\u});
                \draw[bookblue!60,line width=0.7pt,opacity=0.75]
                    plot[domain=-3.6:3.6,samples=55,smooth,variable=\v]
                    ({0.55*\v},{-0.35*\v-0.11*\v*\v});
            \end{scope}
            \draw[bookteal!45,line width=0.65pt]
                ([xshift=3.0cm,yshift=-8.2cm]current page.north west)
                .. controls +(2.2cm,2.0cm) and +(-2.0cm,-2.1cm) ..
                ([xshift=-3.0cm,yshift=9.0cm]current page.south east);
            \draw[bookteal!25,line width=0.55pt]
                ([xshift=3.1cm,yshift=8.0cm]current page.south west)
                .. controls +(2.7cm,-1.7cm) and +(-2.6cm,1.8cm) ..
                ([xshift=-3.1cm,yshift=-8.7cm]current page.north east);
            \node[anchor=west,text=bookteal!72,scale=0.78] at
                ([xshift=2.8cm,yshift=-4.0cm]current page.north west) {%
                $\displaystyle e^{\mathrm{i}\pi}+1=0$};
            \node[anchor=east,text=bookblue!62,scale=0.67] at
                ([xshift=-2.7cm,yshift=4.0cm]current page.south east) {%
                $\displaystyle
                f(x)=\sum_{n=0}^{\infty}\frac{f^{(n)}(a)}{n!}(x-a)^n$};
        \end{tikzpicture}
        \mbox{}
    \end{titlepage}
    \endgroup
}
\fi

% ==================== 新版封面(2026-08 重设计) ====================
% 设计要点:顶部深绿纵向渐变+波浪下缘+两层半透明过渡波(解决深浅绿生硬过渡);
% 深色区螺旋线/同心圆母题;标题+英文副题+作者居中;
% 中部 Fourier 部分和逼近方波曲线;底部浅色波纹(无文字)。
\renewcommand{\maketitle}{%
    \begin{titlepage}
        \thispagestyle{empty}
        \setcounter{page}{0}
        \noindent\hspace*{-\parindent}\begin{tikzpicture}[remember picture,overlay]
            % ---------- 底色 ----------
            \fill[booklight] (current page.north west) rectangle (current page.south east);

            % ---------- 顶部深色区:纵向渐变 + 波浪下缘 ----------
            \shade[top color=bookblue, bottom color=bookteal!55!bookblue]
                (current page.north west) -- (current page.north east) --
                ([yshift=-8.6cm]current page.north east)
                .. controls ([xshift=-5.2cm,yshift=-10.6cm]current page.north east)
                         and ([xshift=6.0cm,yshift=-8.0cm]current page.north west) ..
                ([yshift=-10.0cm]current page.north west) -- cycle;

            % ---------- 过渡层 1:半透明波浪 ----------
            \fill[bookteal!60!bookblue,opacity=0.42]
                ([yshift=-8.6cm]current page.north east)
                .. controls ([xshift=-5.2cm,yshift=-10.6cm]current page.north east)
                         and ([xshift=6.0cm,yshift=-8.0cm]current page.north west) ..
                ([yshift=-10.0cm]current page.north west) --
                ([yshift=-12.6cm]current page.north west)
                .. controls ([xshift=6.5cm,yshift=-11.0cm]current page.north west)
                         and ([xshift=-5.5cm,yshift=-13.2cm]current page.north east) ..
                ([yshift=-11.4cm]current page.north east) -- cycle;

            % ---------- 过渡层 2:更浅的波浪 ----------
            \fill[bookteal!35,opacity=0.35]
                ([yshift=-10.6cm]current page.north east)
                .. controls ([xshift=-6.0cm,yshift=-12.4cm]current page.north east)
                         and ([xshift=5.0cm,yshift=-9.8cm]current page.north west) ..
                ([yshift=-11.8cm]current page.north west) --
                ([yshift=-13.8cm]current page.north west)
                .. controls ([xshift=5.5cm,yshift=-12.6cm]current page.north west)
                         and ([xshift=-6.5cm,yshift=-14.6cm]current page.north east) ..
                ([yshift=-12.8cm]current page.north east) -- cycle;

            % ---------- 深色区装饰:右上螺旋线 ----------
            \begin{scope}[shift={([xshift=-3.4cm,yshift=-4.3cm]current page.north east)}]
                \draw[booklight!45!bookteal,opacity=0.55,line width=0.6pt]
                    plot[domain=0:1440,samples=220,smooth,variable=\t]
                    ({0.0028*\t*cos(\t)},{0.0028*\t*sin(\t)});
            \end{scope}

            % ---------- 深色区装饰:左下同心圆 ----------
            \foreach \r in {0.7,1.25,1.8,2.35,2.9}{
                \draw[booklight!35!bookteal,opacity=0.5,line width=0.5pt]
                    ([xshift=0.4cm,yshift=-7.9cm]current page.north west) circle[radius=\r cm];
            }

            % ---------- 主标题 ----------
            \node[anchor=center,text=booklight] at
                ([yshift=-4.5cm]current page.north)
                {\sffamily\bfseries\fontsize{40}{46}\selectfont \@title};

            % ---------- 分隔短线 ----------
            \draw[booklight!65,line width=1.2pt]
                ([xshift=-2.6cm,yshift=-5.75cm]current page.north) --
                ([xshift=2.6cm,yshift=-5.75cm]current page.north);

            % ---------- 英文副题 ----------
            \node[anchor=center,text=booklight!70!bookteal] at
                ([yshift=-6.55cm]current page.north)
                {\itshape\large Mathematical Analysis \textperiodcentered\
                 A Collection of Exercises};

            % ---------- 作者(紧随标题下方) ----------
            \node[anchor=center,text=booklight!88] at
                ([yshift=-8.2cm]current page.north)
                {\sffamily\large \@author\quad 整理};

            % ---------- 中部:Fourier 部分和逼近方波 ----------
            \begin{scope}[shift={([xshift=-5.2cm,yshift=-18.1cm]current page.north)}]
                % 坐标轴
                \draw[bookgray!55,line width=0.5pt,->] (-0.5,0) -- (10.9,0);
                \draw[bookgray!55,line width=0.5pt,->] (0,-1.75) -- (0,1.95);
                % 目标方波(虚线)
                \draw[bookgray!75,line width=0.6pt,dashed]
                    (0,0.92) -- (5,0.92) (5,-0.92) -- (10,-0.92);
                % N=1
                \draw[bookteal!45,line width=0.7pt]
                    plot[domain=0:360,samples=90,smooth,variable=\x]
                    ({\x/36},{1.18*sin(\x)});
                % N=3
                \draw[bookteal!80,line width=0.8pt]
                    plot[domain=0:360,samples=120,smooth,variable=\x]
                    ({\x/36},{1.18*(sin(\x)+sin(3*\x)/3)});
                % N=7
                \draw[bookblue,line width=0.95pt]
                    plot[domain=0:360,samples=150,smooth,variable=\x]
                    ({\x/36},{1.18*(sin(\x)+sin(3*\x)/3+sin(5*\x)/5+sin(7*\x)/7)});
                % 刻度标注
                \node[anchor=north east,text=bookgray!80,scale=0.62] at (-0.12,0) {$0$};
                \node[anchor=north,text=bookgray!80,scale=0.62] at (5,0) {$\pi$};
                \node[anchor=north,text=bookgray!80,scale=0.62] at (10,0) {$2\pi$};
            \end{scope}
            % 图注
            \node[anchor=center,text=bookgray,scale=0.8] at
                ([yshift=-20.5cm]current page.north)
                {$S_N(x)=\displaystyle\sum_{k=1}^{N}\frac{\sin\bigl((2k-1)x\bigr)}{2k-1}$
                 \ 逐步逼近方波};

            % ---------- 底部浅色波纹(无文字,与顶部波浪呼应) ----------
            \fill[bookteal!22,opacity=0.5]
                ([yshift=2.4cm]current page.south west)
                .. controls ([xshift=6.0cm,yshift=3.9cm]current page.south west)
                         and ([xshift=-6.5cm,yshift=1.0cm]current page.south east) ..
                ([yshift=2.9cm]current page.south east) --
                (current page.south east) -- (current page.south west) -- cycle;
            \fill[bookteal!40!bookblue,opacity=0.28]
                ([yshift=1.3cm]current page.south west)
                .. controls ([xshift=5.5cm,yshift=2.7cm]current page.south west)
                         and ([xshift=-5.5cm,yshift=0.3cm]current page.south east) ..
                ([yshift=1.9cm]current page.south east) --
                (current page.south east) -- (current page.south west) -- cycle;
        \end{tikzpicture}
        \mbox{}
    \end{titlepage}
}

% ==================== 新版封底(纯背景,2026-08 重设计) ====================
% 设计要点:对角渐变底+角部半透明叠层;中下部外摆线(玫瑰形)为视觉主体;
% 底部两层半透明白色波浪;无任何文字。
\newcommand{\makebackcover}{%
    \begingroup
    \let\cleardoublepage\clearpage
    \begin{titlepage}
        \thispagestyle{empty}
        \noindent\hspace*{-\parindent}\begin{tikzpicture}[remember picture,overlay]
            % ---------- 底色:对角渐变 ----------
            \shade[top color=bookblue, bottom color=bookteal!48!bookblue,
                   shading angle=-30]
                (current page.north west) rectangle (current page.south east);

            % ---------- 角部半透明叠层 ----------
            \fill[white,opacity=0.05]
                ([xshift=2.0cm,yshift=1.0cm]current page.north east) circle[radius=7.2cm];
            \fill[white,opacity=0.06]
                ([xshift=-0.5cm,yshift=-1.5cm]current page.north east) circle[radius=3.6cm];

            % ---------- 左下同心圆 ----------
            \foreach \r in {1.0,1.75,2.5,3.25,4.0,4.75}{
                \draw[booklight!40,opacity=0.55,line width=0.55pt]
                    ([xshift=-0.6cm,yshift=-0.4cm]current page.south west) circle[radius=\r cm];
            }
            % ---------- 右下螺旋线 ----------
            \begin{scope}[shift={([xshift=-2.9cm,yshift=3.2cm]current page.south east)}]
                \draw[booklight!45!bookteal,opacity=0.5,line width=0.55pt]
                    plot[domain=0:1260,samples=200,smooth,variable=\t]
                    ({0.0036*\t*cos(\t)},{0.0036*\t*sin(\t)});
            \end{scope}
            % ---------- 左上斜线组 ----------
            \foreach \i in {0,1,2}{
                \draw[bookteal!60!booklight,opacity=0.5,line width=0.6pt]
                    ([xshift={1.6cm+\i*0.45cm}]current page.north west) --
                    ([yshift={-1.6cm-\i*0.45cm}]current page.north west);
            }

            % ---------- 中下部视觉主体:外摆线(玫瑰形) ----------
            \begin{scope}[shift={([yshift=-16.6cm]current page.north)}]
                % 外圈
                \draw[booklight!35,opacity=0.6,line width=0.6pt]
                    (0,0) circle[radius=4.9cm];
                \draw[booklight!22,opacity=0.5,line width=0.5pt]
                    (0,0) circle[radius=5.25cm];
                % 主玫瑰线
                \draw[booklight!55!bookteal,opacity=0.8,line width=0.85pt]
                    plot[domain=0:1080,samples=400,smooth,variable=\t]
                    ({0.335*(8*cos(\t)-5*cos(2.6667*\t))},
                     {0.335*(8*sin(\t)-5*sin(2.6667*\t))});
                % 内层玫瑰线
                \draw[booklight!38!bookteal,opacity=0.6,line width=0.6pt]
                    plot[domain=0:1080,samples=400,smooth,variable=\t]
                    ({0.21*(8*cos(\t)-5*cos(2.6667*\t))},
                     {0.21*(8*sin(\t)-5*sin(2.6667*\t))});
            \end{scope}

            % ---------- 底部半透明波浪 ----------
            \fill[white,opacity=0.06]
                ([yshift=4.6cm]current page.south west)
                .. controls ([xshift=6.5cm,yshift=6.4cm]current page.south west)
                         and ([xshift=-6.0cm,yshift=2.6cm]current page.south east) ..
                ([yshift=5.2cm]current page.south east) --
                (current page.south east) -- (current page.south west) -- cycle;
            \fill[booklight,opacity=0.08]
                ([yshift=2.6cm]current page.south west)
                .. controls ([xshift=5.5cm,yshift=4.3cm]current page.south west)
                         and ([xshift=-6.5cm,yshift=1.0cm]current page.south east) ..
                ([yshift=3.3cm]current page.south east) --
                (current page.south east) -- (current page.south west) -- cycle;
        \end{tikzpicture}
        \mbox{}
    \end{titlepage}
    \endgroup
}
\makeatother

% 页脚右下方常态化水印(署名 + 仓库地址,加粗黑色),已在 main/plain 分页样式中定义,所有页面均生效.

\begin{document}
\mainmatter
\pagestyle{main}
\setcounter{chapter}{7}
\def\BookEnd{\end{document}}
\fi

% ==================== 附录:数学分析重要定理 ====================
% 本附录按小节组织,小节编号 §1,§2,§3,... 自动生成,
% 你只需写 \section{小节名},例如 \section{数列}.
% 子标题自动居左,编号前带段落符号 \S.
%
% 定理请放在 ti 环境(绿色框)中,证明写在框外下方,与正文题目结构一致:
%   \begin{ti}
%   【定理1】 定理陈述...
%   \end{ti}
%   \textbf{证明}:证明内容...
%
% 同样支持 \textbf{另解}、\textbf{注}、\textbf{注记}、\textbf{一点思考} 等标记.
% ================================================================

% 子标题编号格式:§1,§2,§3,...(带段落符号);标题居左.
\renewcommand{\thesection}{\S\arabic{section}}
\setcounter{section}{0}
\titleformat{\section}[block]
  {\normalfont\heiti\bfseries\Large\color{bookblue}}
  {\thesection}{1em}{}
\titlespacing*{\section}{0pt}{3.2em plus 0.4em}{1em}

% 附录内局部放大绿框(ti)之间的间距:+20%(14pt → 16.8pt),不影响正文各章.
\renewtcolorbox{ti}{
    enhanced,
    breakable,
    colback=booklight!55!white,
    colframe=bookteal!35!white,
    boxrule=0.45pt,
    boxsep=0pt,
    left=10pt,
    right=10pt,
    top=9pt,
    bottom=9pt,
    before skip=16.8pt,
    after skip=16.8pt,
    arc=2pt,
    outer arc=2pt,
    before upper={\parindent=0pt}
}

\chapter*{附录\quad 数学分析重要定理}
\phantomsection
\addcontentsline{toc}{chapter}{附录\quad 数学分析重要定理}
\markboth{附录\quad 数学分析重要定理}{}

% 子标题与内容示例(由作者填写),例如:
% \section{数列}
% \begin{ti}
% 【定理1】 单调有界数列必收敛.
% \end{ti}
% \textbf{证明}: ...
%
% \section{级数}
% \begin{ti}
% 【定理1】 正项级数收敛的Cauchy判别法:...
% \end{ti}
% \textbf{证明}: ...

\section{数列}
\begin{ti}
\textbf{定义1.1(数列极限)}  设$\{a_n\}$为数列,$\alpha\in\mathbb{R}$,如果任给$\varepsilon>0$,均存在正整数$N=N(\varepsilon)$,
使得当$n>N$时,$\left\lvert a_n-\alpha\right\rvert<\varepsilon$总成立,则称$\{a_n\}$收敛于$\alpha$,或称$\{a_n\}$以$\alpha$为极限,记为
$$\displaystyle\lim_{n\to\infty}a_n=\alpha\quad\text{或}\quad a_n\rightarrow\alpha$$.
\end{ti}

\begin{ti}
\textbf{命题1.2(绝对值性质与保序性质)} 设数列$\{a_n\}$收敛到$\alpha$,设数列$\{b_n\}$收敛到$\beta$,\\

(1) 则$\{\left\lvert a_n\right\rvert \}$收敛到$\left\lvert \alpha\right\rvert $. \\
(2) 如果存在$N_0$,当$n>N_0$时$a_n\geq b_n$,则$\alpha\geq\beta$. \\
(3) 反之如果$\alpha>\beta$,则存在$N$,使得当$n>N$时$a_n>b_n$.
\end{ti}

\begin{ti}
\textbf{命题1.3(极限的四则运算性质)} 设数列$\{a_n\}$收敛到$\alpha$,设数列$\{b_n\}$收敛到$\beta$,\\

(1)$\{\lambda a_n+\mu b_n\}$收敛到$\lambda\alpha+\mu\beta$,其中$\lambda,\mu$为常数.\\
(2)$\{a_nb_n\}$收敛到$\alpha\beta$.\\
(3)当$\beta\neq 0$时,$\left\{\dfrac{a_n}{b_n}\right\}$收敛到$\dfrac{\alpha}{\beta}$.
\end{ti}

\begin{ti}
\textbf{定理1.4(单调数列的极限)}  设$\{a_n\}$为单调数列. \\

(1)如果$\{a_n\}$为单调递增数列,则$\displaystyle\lim_{n\to\infty}a_n=\sup\left\{a_k\;|\;k\geq 1\right\}$;\\
(2)如果$\{a_n\}$为单调递减数列,则$\displaystyle\lim_{n\to\infty}a_n=\inf\left\{a_k\;|\;k\geq 1\right\}$.
\end{ti}

\begin{ti}
\textbf{推论1.5} 有界单调数列必收敛
\end{ti}

\begin{ti}
\textbf{定义1.6(上下极限)} 对于有界数列$\{a_n\}$,令
$$\underline{a_n}=\inf\left\{a_k\;|\;k\geq n\right\},\quad \overline{a_n}=\sup\left\{a_k\;|\;k\geq n\right\}$$.
易见$\{\underline{a_n}\}$和$\{\overline{a_n}\}$分别是单调递增和单调递减的数列,且
$$\underline{a_n}\leq a_n\leq\overline{a_n}$$.
单调数列$\{\underline{a_n}\}$和$\{\overline{a_n}\}$的极限称为$\{a_n\}$的\textbf{下极限}和\textbf{上极限},记为
$$\displaystyle\varliminf_{n\to\infty}a_n=\lim_{n\to\infty}\underline{a_n},\quad\varlimsup_{n\to\infty}a_n=\lim_{n\to\infty}\overline{a_n}$$.
\end{ti}

\begin{ti}
\textbf{定理1.7} 设$\{a_n\}$为有界数列,则下列命题等价: \\

(1)$\{a_n\}$收敛;\\
(2)$\{a_n\}$的上极限和下极限相等;\\
(3)$\displaystyle\lim_{n\to\infty}(\overline{a_n}-\underline{a_n})=0$.
\end{ti}

\begin{ti}
\textbf{命题1.8} 设$\{a_n\},\{b_n\}$为有界数列.\\
(1)(上下极限的保序性质)如果存在$N_0,n>N_0$时$a_n\geq b_n$,则
$$\displaystyle\varliminf_{n\to\infty}a_n\geq\varliminf_{n\to\infty}b_n,\quad\varlimsup_{n\to\infty}a_n\geq\varlimsup_{n\to\infty}b_n$$;
(2)$\displaystyle\varlimsup_{n\to\infty}(a_n+b_n)\leq\varlimsup_{n\to\infty}a_n+\varlimsup_{n\to\infty}b_n$.
\end{ti}

\begin{ti}
\textbf{定理1.9(Cauchy准则)} 数列$\{a_n\}$收敛当且仅当它是$Cauchy$数列.
\end{ti}

\textbf{证明}:必要性:设$\{a_n\}$收敛到$\alpha$.任给$\varepsilon>0$,存在$N$,当$n>N$时$\left\lvert a_n-\alpha\right\rvert<\dfrac{\varepsilon}{2}$.因此,当$m,n>N$时,有
$$\left\lvert a_m-a_n\right\rvert\leq\left\lvert a_m-\alpha\right\rvert+\left\lvert\alpha-a_n\right\rvert<\frac{\varepsilon}{2}+\frac{\varepsilon}{2}=\varepsilon,$$
这说明$\{a_n\}$为$Cauchy$数列. \\

\phantom{\textbf{证明}:}充分性:设$\{a_n\}$为$Cauchy$数列.由于$Cauchy$数列必有界,$\{a_n\}$是有界数列,于是可以研究其上、下极限.根据$Cauchy$数列的定义,
任给$\varepsilon>0$,存在$N$,当$m,n>N$时,$\left\lvert a_m-a_n\right\rvert<\varepsilon$.
在上式中暂时固定$n>N$,对$\{a_m\}$取上极限,利用上极限的保序性可得$-\varepsilon\leq\displaystyle\varlimsup_{m\to\infty}a_m-a_n\leq\varepsilon$.
由数列极限的定义即可看出$\{a_n\}$收敛.

\begin{ti}
\textbf{定理1.10(Stolz公式)} 设$\{x_n\},\{y_n\}$为数列, \\

(1)若$\{y_n\}$严格单调地趋于$+\infty,\displaystyle\lim_{n\to\infty}\dfrac{x_n-x_{n-1}}{y_n-y_{n-1}}=\alpha,\alpha$为实数或$\pm\infty$,则
$$\displaystyle\lim_{n\to\infty}\dfrac{x_n}{y_n}=\alpha$$. \\
(2)若$\{y_n\}$严格单调递减趋于$0,\{x_n\}$也收敛到$0,\displaystyle\lim_{n\to\infty}\dfrac{x_n-x_{n-1}}{y_n-y_{n-1}}=\alpha,\alpha$为实数或$\pm\infty$,则
$$\displaystyle\lim_{n\to\infty}\dfrac{x_n}{y_n}=\alpha$$.
\end{ti}

\textbf{证明}:先证(1).分情况讨论. \\
\romannum{1}$\;\alpha\in\mathbb{R}$.先设$\alpha=0$.此时,任给$\varepsilon>0$,存在$N$,当$n>N$时,$\left\lvert\dfrac{x_n-x_{n-1}}{y_n-y_{n-1}}\right\rvert<\varepsilon$.
利用$\{y_n\}$的单调性以及引理(当$1\leq k\leq n$时,设$b_k>0$且$m\leq\dfrac{a_k}{b_k}\leq M$,则$m\leq\dfrac{a_1+\cdots+a_n}{b_1+\cdots+b_n}\leq M$),当$n>N$时,得到
$$\left\lvert\dfrac{x_n-x_N}{y_n-y_N}\right\rvert<\varepsilon,$$
整理以后可得
$$\frac{x_N-\varepsilon y_N}{y_n}<\frac{x_n}{y_n}<\frac{x_N+\varepsilon y_N}{y_n}+\varepsilon.$$
由$\displaystyle\lim_{n\to\infty}y_n=+\infty$以及上式可知,当$n$充分大时,$-2\varepsilon<\dfrac{x_n}{y_n}<2\varepsilon$,因此$\displaystyle\lim_{n\to\infty}\dfrac{x_n}{y_n}=0$. \\
一般地,设$\displaystyle\lim_{n\to\infty}\dfrac{x_n-x_{n-1}}{y_n-y_{n-1}}=\alpha$,记$\widetilde{x_n}=x_n-\alpha y_n$,则
$$\lim_{n\to\infty}\dfrac{\widetilde{x_n}-\widetilde{x_{n-1}}}{y_n-y_{n-1}}=\lim_{n\to\infty}\left(\dfrac{x_n-x_{n-1}}{y_n-y_{n-1}}-\alpha\right)=0,$$
这说明$\displaystyle\lim_{n\to\infty}\dfrac{\widetilde{x_n}}{y_n}=0$,从而$\displaystyle\lim_{n\to\infty}\dfrac{x_n}{y_n}=\alpha$. \\

\romannum{2}$\;\alpha=+\infty$.此时,存在$N$,当$n>N$时,$\dfrac{x_n-x_{n-1}}{y_n-y_{n-1}}>1$.于是,当$n>N$时,$x_n-x_{n-1}>y_n-y_{n-1}>0$,即$n>N$时$\{x_n\}$也是严格单调递增的,且
$$x_n-x_N=(x_n-x_{n-1})+\cdots+(x_{N+1}-x_N)>(y_n-y_N),$$
由$\displaystyle\lim_{n\to\infty}y_n=+\infty$可知$\displaystyle\lim_{n\to\infty}x_n=+\infty$.
由题设可得$\displaystyle\lim_{n\to\infty}\dfrac{y_n-y_{n-1}}{x_n-x_{n-1}}=0$.
应用(1)的结论可得$\displaystyle\lim_{n\to\infty}\dfrac{y_n}{x_n}=0$,于是$\displaystyle\lim_{n\to\infty}\dfrac{x_n}{y_n}=+\infty$. \\

\romannum{3}$\;\alpha=-\infty$.这时只要将$x_n$换成$-x_n$,然后应用(2)的结论即可. \\

再证(2). \\
\romannum{1}$\;\alpha\in\mathbb{R}$.不妨设$\alpha=0$.任给$\varepsilon>0$,存在$N$,当$n>N$时,$\left\lvert\dfrac{x_{n+1}-x_n}{y_{n+1}-y_n}\right\rvert<\varepsilon$,和上面的证明类似,当$m>n>N$时可得
$$-\varepsilon(y_n-y_m)<x_n-x_m<\varepsilon(y_n-y_m),$$
令$m\to\infty$,可得$-\varepsilon y_n\leq x_n\leq\varepsilon y_n$,这说明$\displaystyle\lim_{n\to\infty}\dfrac{x_n}{y_n}=0$. \\

\romannum{2}$\;\alpha=+\infty$.任给$M>0$,存在$N$,当$n>N$时,$\dfrac{x_n-x_{n+1}}{y_n-y_{n+1}}>M$.
与上面的证明类似,当$m>n>N$时,有$x_n-x_m>M(y_n-y_m)$,令$m\to\infty$,得$x_n\geq My_n$,这说明$\displaystyle\lim_{n\to\infty}\dfrac{x_n}{y_n}=+\infty$. \\

\romannum{3}$\;\alpha=-\infty$.将(2)中$x_n$换成$-x_n$即可.

\section{实数系基本定理}
\begin{ti}
\textbf{定理2.1(Cantor)} 设$\{[a_n,b_n]\}$是闭区间套,如果$\displaystyle\lim_{n\to\infty}(b_n-a_n)=0,$则这些闭区间有一个唯一的公共点.
\end{ti}

\textbf{证明}:由题设可知,$\{a_n\}$单调递增,$\{b_n\}$单调递减,它们均位于区间$[a_1,b_1]$中,从而为有界数列.
由单调有界数列必收敛,$\{a_n\}$和$\{b_n\}$都收敛,设其极限分别为$\alpha,\beta$.
由$a_n<b_n$和极限的保序性质可知$\alpha\leq\beta$.由$\{a_n\}$和$\{b_n\}$的单调性可知$a_n\leq\alpha,\beta\leq b_n$.
于是$\alpha,\beta$为$\{[a_n,b_n]\}$的公共点,且
$$0\leq\beta-\alpha\leq b_n-a_n,\;\forall\,n\geq 1.$$
令$n\to\infty$,由$\displaystyle\lim_{n\to\infty}(b_n-a_n)=0$即得$\alpha=\beta$.
如果$\{[a_n,b_n]\}$另有公共点$\gamma$,则
$$0\leq\left\lvert\gamma-\alpha\right\rvert\leq b_n-a_n,\;\forall\,n\geq 1,$$同理可得$\gamma=\alpha$.

\begin{ti}
\textbf{定理2.2(Heine-Borel)} 设$\{U_{\alpha}\}_{\alpha\in\varGamma}$为$\mathbb{R}$中的一族开集,如果闭区间$[a,b]$包含于这一族开集的并集之中,则$[a,b]$
必定包含于有限个$U_{\alpha}$的并集之中.
\end{ti}

\textbf{证明}:如果指标集$\varGamma$是有限集,则结论不证自明.以下设$\varGamma$为无限集.\\
(反证法)假设$[a,b]$只能包含于无限个$U_{\alpha}$的并集,则二等分$[a,b]$后必有一个小区间也只能包含于无限个$U_{\alpha}$之并,记该区间为$[a_1,b_1]$.
再将$[a_1,b_1]$二等分,又必有一个小区间只能包含于无限个$U_{\alpha}$之并,记该区间为$[a_2,b_2]$.
如此继续下去,得闭区间套$\{[a_n,b_n]\}$,使得每一个$[a_n,b_n]$均只能包含于无限个$U_{\alpha}$之并.\\

注意$\displaystyle\lim_{n\to\infty}(b_n-a_n)=\lim_{n\to\infty}2^{-n}(b-a)=0$.根据闭区间套原理,$\{[a_n,b_n]\}$有一个公共点,记为$\xi$.
根据题设,存在$\alpha_0\in\varGamma$,使得$\xi\in U_{\alpha_0}$.因为$U_{\alpha_0}$是开集,故存在$\delta>0$,使得$(\xi-\delta,\xi+\delta)\subset U_{\alpha_0}$.
根据闭区间套原理的证明,$\{a_n\},\{b_n\}$均收敛于$\xi$,故存在$N$,当$n>N$时$a_n,b_n\in(\xi-\delta,\xi+\delta)$.此时有
$$[a_n,b_n]\subset(\xi-\delta,\xi+\delta)\subset U_{\alpha_0},$$
这与$[a_n,b_n]$只能包含于无限个$U_{\alpha}$之并相矛盾.

\begin{ti}
\textbf{定理2.3(Bolzano)}有界数列必有收敛子列
\end{ti}

\textbf{证明}:设$\{a_n\}$为有界数列,不妨设$a_n$均包含于$[a,b]$.\\
\phantom{\textbf{证明}:}断言:存在$\alpha\in[a,b]$,使得任给$\delta>0$,$(\alpha-\delta,\alpha+\delta)$中均含有无限项$a_n$.\\
\phantom{\textbf{证明}:}(反证法)假设不然,则任给$x\in[a,b]$,存在$\delta(x)>0$,使得$(x-\delta(x),x+\delta(x))$只含有限项$a_n$.
显然,$[a,b]$包含于集合族$\{(x-\delta(x),x+\delta(x))\}_{x\in[a,b]}$之并.由$Heine-Borel$定理,$[a,b]$包含于有限个$(x-\delta(x),x+\delta(x))$之并.
这说明$[a,b]$只含有限项$a_n$,与$a_n$均包含于$[a,b]$相矛盾.利用此断言,我们选取$\{a_n\}$的子列,使之收敛到$\alpha$.
事实上,先取$a_{n_1}\in(\alpha-1,\alpha+1)$.再取$n_2>n_1$,使得$a_{n_2}\in\left(\alpha-\dfrac{1}{2},\alpha+\dfrac{1}{2}\right)$.
如此继续,可得子列$\{a_{n_k}\}$,使得当$k\geq 1$时$a_{n_k}\in\left(\alpha-\dfrac{1}{k},\alpha+\dfrac{1}{k}\right)$.显然,$\{a_{n_k}\}$收敛到$\alpha$.

\section{函数}
\begin{ti}
\textbf{定义3.1} 设函数$f$在$x_0$的一个空心开邻域$V(x_0,\delta_0)$中有定义.如果存在$\alpha\in\mathbb{R}$,使得任给$\varepsilon>0$,
均存在$\delta\in(0,\delta_0)$,当$x\in V(x_0,\delta)$时,$\left\lvert f(x)-\alpha\right\rvert<\varepsilon$,则称函数$f$在$x_0$处有极限$\alpha$,记为
$$\displaystyle\lim_{x\to x_0}f(x)=\alpha\;\text{或}\;f(x)\rightarrow\alpha(x\rightarrow x_0)$$.
\end{ti}

\begin{ti}
\textbf{命题3.2} \\
(1)(局部有界性)如果$f$在$x_0$处有极限,则它在$x_0$的一个空心开邻域中有界.\\
(2)(保序性)设函数$f,g$在$x_0$处的极限分别为$\alpha,\beta$.如果$f(x)\geq g(x)$在$x_0$的一个空心开邻域中成立,则$\alpha\geq\beta$,
反之,如果$\alpha>\beta$,则在$x_0$的一个空心开邻域中$f(x)>g(x)$;特别地,如果$\alpha>0$,则在$x_0$的一个空心开邻域中$f(x)>0$
\end{ti}

\begin{ti}
\textbf{定理3.3(复合函数的极限)}设函数$f(y)$在$y_0$处极限为$\alpha$,函数$g(x)$在$x_0$处极限为$y_0$,
如果存在$x_0$的一个空心开邻域,在此开邻域中$g(x)\neq y_0$,则复合函数$f(g(x))$在$x_0$处极限为$\alpha$
\end{ti}

\textbf{证明}:任给$\varepsilon>0$,由$\displaystyle\lim_{y\to y_0}f(y)=\alpha$可知,存在$\delta>0$,当$0<\left\lvert y-y_0\right\rvert<\delta$时,
$\left\lvert f(y)-\alpha\right\rvert<\varepsilon$.因为$\displaystyle\lim_{x\to x_0}g(x)=y_0$,
对于这个$\delta$,存在$\eta>0$,使得当$0<\left\lvert x-x_0\right\rvert<\eta$时,$0<\left\lvert g(x)-y_0\right\rvert<\delta$.
此时$\left\lvert f(g(x))-\alpha\right\rvert<\varepsilon$,这说明$f(g(x))$在$x_0$处的极限为$\alpha$.

\begin{ti}
\textbf{定理3.4(Heine)} 设函数$f$在$x_0$的一个空心开邻域中有定义,则$f$在$x_0$
处极限为$\alpha$当且仅当对空心开邻域中任何收敛于$x_0$的数列$\{x_n\}$,均有
$\displaystyle\lim_{n\to\infty}f(x_n)=\alpha$.
\end{ti}

\textbf{证明}:必要性:由函数极限的定义,任给$\varepsilon>0$,存在$\delta>0$,使得当$x\in V(x_0,\delta)$时,$\left\lvert f(x)-\alpha\right\rvert<\varepsilon$.
现在设$\{x_n\}$是空心开邻域中收敛于$x_0$的数列,由数列极限的定义,对于上述$\delta$,存在$N$,使得当$n>N$时$0<\left\lvert x_n-x_0\right\rvert<\delta$.
这说明当$n>N$时,$\left\lvert f(x_n)-\alpha\right\rvert<\varepsilon$,因此$\displaystyle\lim_{n\to\infty}f(x_n)=\alpha$. \\
\phantom{\textbf{证明}:}充分性:(反证法)此时,根据函数极限的定义,存在$\varepsilon_0>0$,使得任给$n\geq 1$,均存在$x_n\in V\left(x_0,\dfrac{1}{n}\right)$,使得$\left\lvert f(x_n)-\alpha\right\rvert\geq\varepsilon_0$.
这说明$\{x_n\}$是空心开邻域中收敛于$x_0$的数列,但$\{f(x_n)\}$不收敛到$\alpha$.

\begin{ti}
\textbf{定义3.5(无穷小量与无穷大量)}  设$x_0$为实数或$\pm\infty$.\\
(1)如果$\displaystyle\lim_{x\to x_0}f(x)=0$,则称函数$f$在$x\to x_0$时为\textbf{无穷小量},记为$f(x)=o(1)\,(x\to x_0)$;\\
(2)如果$\dfrac{1}{f}$在$x\to x_0$时为无穷小量,则称$f$在$x\to x_0$时为\textbf{无穷大量}.
\end{ti}

\begin{ti}
\textbf{命题3.6(无穷小量的比较)}  设$f,g$均为$x\to x_0$时的无穷小量.\\
(1)如果$\dfrac{f(x)}{g(x)}=o(1)\,(x\to x_0)$,则称当$x\to x_0$时$f$关于$g$为\textbf{高(低)阶}无穷小量,记为$f(x)=o(g(x))\,(x\to x_0)$;\\
(2)如果$\dfrac{f(x)}{g(x)}$在$x_0$处的极限为非零实数,则称当$x\to x_0$时$f$与$g$为\textbf{同阶量};\\
(3)如果$\dfrac{f(x)}{g(x)}$在$x_0$处的极限等于$1$,则称当$x\to x_0$时$f$与$g$为\textbf{等价量},记为$f\sim g\,(x\to x_0)$.
\end{ti}

\begin{ti}
\textbf{定义3.7(函数的连续性)}  设函数$f$在$x_0$的某个邻域中有定义.如果$\displaystyle\lim_{x\to x_0}f(x)=f(x_0)$,则称$f$在$x_0$处\textbf{连续}.
如果$f$在区间$I$的每一点处均连续,则称$f$在$I$上连续.
\end{ti}

\begin{ti}
\textbf{命题3.8(连续函数的运算性质)}  设函数$f,g$在$x_0$处连续.\\
(1)$f\pm g$, $fg$以及$cf$($c$为常数)均在$x_0$处连续;若$g(x_0)\neq 0$,则$\dfrac{f}{g}$也在$x_0$处连续;\\
(2)设$f$在$x_0$处连续,$g$在$y_0=f(x_0)$处连续,则复合函数$g\circ f$在$x_0$处连续.
\end{ti}



\begin{ti}
\textbf{定义3.9(间断点)}  设函数$f$在$x_0$的某个空心开邻域中有定义.如果$f$在$x_0$处不连续,则称$x_0$为$f$的\textbf{间断点}.
若$f$在$x_0$处的左右极限均存在,则称$x_0$为\textbf{第一类间断点},其中左右极限不等者称为\textbf{跳跃间断点},左右极限相等者称为\textbf{可去间断点};
其余情形称为\textbf{第二类间断点}.
\end{ti}

\begin{ti}
\textbf{定理3.10(最值定理)}  设函数$f$在闭区间$[a,b]$上连续,则$f$在$[a,b]$中必定取到最大值和最小值,即存在$x_1,x_2\in[a,b]$,使得
$$f(x_1)\leq f(x)\leq f(x_2),\quad \forall x\in[a,b].$$
\end{ti}

\textbf{证明}:由连续函数的有界性质,$f$有界,因此必有上确界和下确界.记上确界为$M$.
根据确界的刻画,存在$[a,b]$中的数列$\{x_n\}$,使得$M-\frac{1}{n}\leq f(x_n)\leq M$.根据$Bolzano$定理,$\{x_n\}$有收敛子列$\{x_{n_i}\}$,其极限记为$x_*$.
由极限的保序性质可知$x_*\in[a,b]$.因为$f$在$x_*$处连续,故$\{f(x_{n_i})\}$收敛于$f(x_*)$.这说明$f(x_*)=M$,$M$即为$f$的最大值.同理可证$f$可以取到最小值(或考虑$-f$的最大值).

\begin{ti}
\textbf{定理3.11(介值定理)}  设函数$f$在闭区间$[a,b]$上连续,$\mu$是严格介于$f(a)$与$f(b)$之间的数,则存在$\xi\in(a,b)$,使得$f(\xi)=\mu$.
\end{ti}

\textbf{证明}:设$\mu$是严格介于$f(a)$和$f(b)$之间的数,则$(f(a)-\mu)(f(b)-\mu)<0$.因此,由零值定理,连续函数$f(x)-\mu$在$(a,b)$内存在零点$\xi$,此时$f(\xi)=\mu$.

\begin{ti}
\textbf{推论3.12(零值定理,Bolzano)}  设函数$f$在闭区间$[a,b]$上连续,且$f(a)f(b)\leq 0$,则存在$\xi\in[a,b]$,使得$f(\xi)=0$.
\end{ti}

\begin{ti}
\textbf{定义3.13(一致连续)}  设函数$f$在区间$I$上有定义.如果对任意的$\varepsilon>0$,均存在$\delta>0$,使得对任意的$x_1,x_2\in I$,
只要$\left\lvert x_1-x_2\right\rvert<\delta$,就有$\left\lvert f(x_1)-f(x_2)\right\rvert<\varepsilon$,则称$f$在$I$上\textbf{一致连续}.
\end{ti}

\begin{ti}
\textbf{定理3.14(Cantor)}  设函数$f$在闭区间$[a,b]$上连续,则$f$在$[a,b]$上一致连续.
\end{ti}

\textbf{证明}:设$f\in C^0[a,b]$.(反证法)若$f$不一致连续,则存在$\varepsilon_0>0$,以及$[a,b]$中的数列$\{a_n\},\{b_n\}$,使得
$$\lim_{n\to\infty}(a_n-b_n)=0,\qquad \left\lvert f(a_n)-f(b_n)\right\rvert\geq\varepsilon_0,\;\forall n\geq 1.$$
\phantom{\textbf{证明}:}根据$Bolzano$定理,$\{b_n\}$有收敛子列$\{b_{n_i}\}$,设其极限为$\xi$,
则$\xi\in[a,b]$.此时$a_{n_i}=(a_{n_i}-b_{n_i})+b_{n_i}\to 0+\xi=\xi\;(i\to\infty)$.因为$f$在$\xi$处连续,故
$$\varepsilon_0\leq\left\lvert f(a_{n_i})-f(b_{n_i})\right\rvert\to\left\lvert f(\xi)-f(\xi)\right\rvert=0\;(i\to\infty),$$
这就导出了矛盾.

\begin{ti}
\textbf{定义3.15(连续函数的定积分)}  \\
设函数$f$在闭区间$[a,b]$上连续.任取$[a,b]$的分割$P$
$$a=x_0<x_1<\cdots<x_n=b$$
以及$\xi_i\in[x_{i-1},x_i]$,作$Riemann$和$\displaystyle\sum_{i=1}^{n}f(\xi_i)\Delta x_i$,其中$\Delta x_i=x_i-x_{i-1}$.
当分割$P\text{的模}\left\lVert P\right\rVert =\displaystyle\max_{1\leq i\leq n}\Delta x_i\to 0$时,上述$Riemann$和的极限称为$f$在$[a,b]$上的\textbf{定积分},记为
$$\int_a^b f(x)\,\mathrm{d}x=\lim_{\left\lVert P\right\rVert\to 0}\sum_{i=1}^{n}f(\xi_i)\Delta x_i$$.
\end{ti}

\begin{ti}
\textbf{命题3.16(定积分的基本性质)}  设函数$f,g$在$[a,b]$上连续.\\
(1)(线性叠加性质)对任意常数$\lambda,\mu$,有
$$\int_a^b\bigl[\lambda f(x)+\mu g(x)\bigr]\,\mathrm{d}x=\lambda\int_a^b f(x)\,\mathrm{d}x+\mu\int_a^b g(x)\,\mathrm{d}x;$$
(2)(区间可加性)对任意$c\in(a,b)$,有
$$\int_a^b f(x)\,\mathrm{d}x=\int_a^c f(x)\,\mathrm{d}x+\int_c^b f(x)\,\mathrm{d}x;$$
(3)(单调性)如果在$[a,b]$上$f(x)\leq g(x)$恒成立,则$\displaystyle\int_a^b f(x)\,\mathrm{d}x\leq\int_a^b g(x)\,\mathrm{d}x$.
\end{ti}

\begin{ti}
\textbf{定理3.17(积分中值定理)}  设函数$f,g$在$[a,b]$上连续,如果$g$不变号,则存在$\xi\in[a,b]$,使得
$$\int_a^b f(x)g(x)\,\mathrm{d}x=f(\xi)\int_{a}^{b} g(x) \,dx .$$
特别地,在上式令$g\equiv 1$,得到$\int_{a}^{b} f(x) \,dx =f(\xi)(b-a)$
\end{ti}

\textbf{证明}:不妨设$g\geq 0$.根据最值定理,$f$在$[a,b]$中有最小值和最大值,分别记为$\alpha,\beta$,则$\alpha g\leq fg\leq\beta g$.根据积分的保序性质和线性叠加性质,有
$$\alpha\int_a^b g(x)\,\mathrm{d}x\leq\int_a^b f(x)g(x)\,\mathrm{d}x\leq\beta\int_a^b g(x)\,\mathrm{d}x.$$
如果$g$在$[a,b]$中的积分为零,则上式表明$fg$的积分也为零,我们可以任意选取$\xi$;如果$g$的积分为正,则
$$\alpha\leq\left(\int_a^b g(x)\,\mathrm{d}x\right)^{-1}\int_a^b f(x)g(x)\,\mathrm{d}x\leq\beta,$$
由介值定理可知,存在$\xi\in[a,b]$,使得$f(\xi)=\left(\displaystyle\int_a^b g(x)\,\mathrm{d}x\right)^{-1}\displaystyle\int_a^b f(x)g(x)\,\mathrm{d}x$.
特别地,当$g\equiv 1$时,上式成为$f(\xi)=\dfrac{1}{b-a}\displaystyle\int_a^b f(x)\,\mathrm{d}x$,即$\displaystyle\int_a^b f(x)\,\mathrm{d}x=f(\xi)(b-a)$.

\section{一元微分学与一元积分学}
\begin{ti}
\textbf{定义4.1(导数)}  设函数$f$在$x_0$附近有定义.如果极限
$$\lim_{x\to x_0}\frac{f(x)-f(x_0)}{x-x_0}$$
存在且有限,则称$f$在$x_0$处\textbf{可导},此极限称为$f$在$x_0$处的\textbf{导数},记为$f'(x_0)$.
\end{ti}


\begin{ti}
\textbf{命题4.2(导数的四则运算)}  设$f,g$在$x_0$处可导,$\lambda,\mu\in\mathbb{R}$,则$fg$和$\lambda f+\mu g$也在$x_0$处可导,且\\
(1)$(\lambda f+\mu g)'(x_0)=\lambda f'(x_0)+\mu g'(x_0)$;\\
(2)$(fg)'(x_0)=f'(x_0)g(x_0)+f(x_0)g'(x_0)$;\\
(3)当$g(x_0)\neq 0$时,$\dfrac{f}{g}$也在$x_0$处可导,且$\left(\dfrac{f}{g}\right)'(x_0)=\dfrac{f'(x_0)g(x_0)-f(x_0)g'(x_0)}{g^2(x_0)}$.
\end{ti}

\begin{ti}
\textbf{定理4.3(链式法则)}  设$g$在$x_0$处可导,$f$在$g(x_0)$处可导,则复合函数$f\circ g=f(g)$在$x_0$处可导,且
$$[f(g)]'(x_0)=f'(g(x_0))g'(x_0)$$.
\end{ti}

\textbf{证明}:记$y_0=g(x_0)$.沿用微分定义中的记号,我们知道$R_f(y,y_0)=o(y-y_0)\,(y\to y_0)$.根据可导函数在$x_0$附近有界,$g(x)-y_0=O(x-x_0)\,(x\to x_0)$.于是利用复合函数的极限可得$R_f(g(x),y_0)=o(x-x_0)\,(x\to x_0)$.这说明
$$f(g(x))=f(y_0)+f'(y_0)(g(x)-y_0)+R_f(g(x),y_0)=f(y_0)+f'(y_0)g'(x_0)(x-x_0)+o(x-x_0),$$
因此$f(g)$在$x_0$处可微(可导),其导数为$f'(y_0)g'(x_0)$.

\begin{ti}
\textbf{定理4.4(反函数求导法则)}  设$f$在$x_0$附近连续且有反函数$g$.如果$f$在$x_0$处可导,且导数$f'(x_0)\neq 0$,则$g$在$y_0=f(x_0)$处可导,且
$$g'(y_0)=\left[f'(x_0)\right]^{-1}$$.
\end{ti}

\begin{ti}
\textbf{定义4.5(微分)}  设函数$f$在$x_0$的某邻域中有定义.如果存在$\alpha\in\mathbb{R}$,使得
$$f(x)=f(x_0)+\alpha(x-x_0)+o(x-x_0)\,(x\rightarrow x_0),$$
则称$f$在$x_0$处\textbf{可微},线性映射
$$df(x_0):\mathbb{R} \rightarrow \mathbb{R} ,t\rightarrow \alpha t$$
称为$f$在$x_0$处的\textbf{微分}.
\end{ti}

\begin{ti}
\textbf{定理4.6(Newton-Leibniz公式)}  设函数$f$在$[a,b]$上连续,且函数$F$为$f$的一个原函数,则
$$\int_a^b f(x)\,\mathrm{d}x=F(b)-F(a)$$.
\end{ti}

\textbf{证明}:记$h(x)=\displaystyle\int_a^x f(t)\,\mathrm{d}t$,由变上限积分定理可知$h'=f$.于是$(F-h)'=0$处处成立,故$F-h$为常值函数.特别地,有
$$\int_a^b f(x)\,\mathrm{d}x=h(b)-h(a)=F(b)-F(a).$$

\begin{ti}
\textbf{命题4.7(分部积分法)}  设函数$u,v$在$[a,b]$上有连续导数,则
$$\int_a^b u(x)v'(x)\,\mathrm{d}x=u(x)v(x)\Big|_a^b-\int_a^b u'(x)v(x)\,\mathrm{d}x$$.
\end{ti}

\begin{ti}
\textbf{命题4.8(换元积分法)}  设$\varphi\in C^1[a,b]$,函数$f$在区间$I\supset\varphi([a,b])$中连续,则
$$\int_a^b f(\varphi(x))\varphi'(x)\,\mathrm{d}x=\int_{\varphi(a)}^{\varphi(b)}f(y)\,\mathrm{d}y$$.
\end{ti}

\begin{ti}
\textbf{定义4.9(极值)}  设函数$f$在$x_0$的某邻域中有定义.如果存在$\delta>0$,使得当$\left\lvert x-x_0\right\rvert<\delta$时$f(x)\leq f(x_0)$恒成立,
则称$x_0$为$f$的\textbf{极大值点},$f(x_0)$为\textbf{极大值};极小值点与极小值类似定义.极大值与极小值统称\textbf{极值}.
\end{ti}

\begin{ti}
\textbf{定理4.10(Fermat引理)}  设$x_0$为函数$f$在$I$中的极值点.如果$f$在$x_0$处可导,且$x_0$为$I$的内点,则$f'(x_0)=0$.
\end{ti}

\textbf{证明}:不妨设$x_0$为$f$的极大值点(不然可考虑$-f$).由于$x_0$为$I$的内点,故存在$\delta>0$,使得$(x_0-\delta,x_0+\delta)\subset I$.取$\delta$充分小,使得$f(x)\leq f(x_0)$对一切$x\in(x_0-\delta,x_0+\delta)$成立.特别地,当$x\in(x_0-\delta,x_0)$时,
$$f'(x_0)=f'_-(x_0)=\lim_{x\to x_0^-}\frac{f(x)-f(x_0)}{x-x_0}\geq 0;$$
\phantom{\textbf{证明}:}而当$x\in(x_0,x_0+\delta)$时,
$$f'(x_0)=f'_+(x_0)=\lim_{x\to x_0^+}\frac{f(x)-f(x_0)}{x-x_0}\leq 0.$$
这说明$f'(x_0)=0$.

\begin{ti}
\textbf{定理4.11(Rolle中值定理)}  设函数$f$在$[a,b]$上连续,在$(a,b)$内可导,且$f(a)=f(b)$,则存在$\xi\in(a,b)$,使得$f'(\xi)=0$.
\end{ti}

\textbf{证明}:根据最值定理,$f$有最大值和最小值.如果二者相等,则$f$为常值函数,定理自然成立;否则由$f(a)=f(b)$可知最大值和最小值中至少有一个是被$f$在内点$\xi\in(a,b)$处所取得.由$Fermat$定理,$f'(\xi)=0$.

\begin{ti}
\textbf{定理4.12(Lagrange中值定理)}  设函数$f$在$[a,b]$上连续,在$(a,b)$内可导,则存在$\xi\in(a,b)$,使得
$$f(b)-f(a)=f'(\xi)(b-a)$$.
\end{ti}

\textbf{证明}:记$\nu=\dfrac{f(b)-f(a)}{b-a}$,$g(x)=f(x)-\nu x$,则$g(b)=g(a)$.根据$Rolle$定理,存在$\xi\in(a,b)$,使得$g'(\xi)=0$.此时$f'(\xi)=\nu$,即$f(b)-f(a)=f'(\xi)(b-a)$.

\begin{ti}
\textbf{定理4.13(Cauchy中值定理)}  设函数$f,g$在$[a,b]$上连续,在$(a,b)$内可导,且$g'(x)\neq 0$,则存在$\xi\in(a,b)$,使得
$$\frac{f(b)-f(a)}{g(b)-g(a)}=\frac{f'(\xi)}{g'(\xi)}$$.
\end{ti}

\textbf{证明}:由$g'$处处非零和$Rolle$定理可知$g(b)\neq g(a)$.记$\nu=\dfrac{f(b)-f(a)}{g(b)-g(a)}$,$h(x)=f(x)-\nu g(x)$,则$h(a)=h(b)$.根据$Rolle$定理,存在$\xi\in(a,b)$,使得$h'(\xi)=0$.此时$f'(\xi)=\nu g'(\xi)$,整理以后即得欲证等式.

\begin{ti}
\textbf{定义4.14(凸函数)}  设函数$f$定义在区间$I$上.如果对任意的$x,y\in I$和任意的$t\in[0,1]$,均有
$$f((1-t)x+ty)\leq(1-t)f(x)+tf(y)$$
成立,则称$f$为$I$上的\textbf{凸函数};若当$x\neq y$且$t\in(0,1)$时不等号严格成立,则称$f$为\textbf{严格凸函数}.
\end{ti}

\begin{ti}
\textbf{定理4.15(Jensen不等式)}  设函数$f$为区间$I$上的凸函数.对任意的$x_1,\ldots,x_n\in I$以及满足$\displaystyle\sum_{i=1}^{n}\lambda_i=1$的非负实数$\lambda_i$,有
$$f\left(\sum_{i=1}^{n}\lambda_i x_i\right)\leq\sum_{i=1}^{n}\lambda_i f(x_i)$$.
\end{ti}

\textbf{证明}:对$n$用数学归纳法.$n=1$是显然的,$n=2$由凸函数定义直接得到.假设不等式对$n=k$成立.当$n=k+1$时,不妨设$0<\lambda_{k+1}<1$,此时$\displaystyle\sum_{i=1}^{k}\frac{\lambda_i}{1-\lambda_{k+1}}=1$.由归纳假设,有
$$f\left(\sum_{i=1}^{k+1}\lambda_i x_i\right)=f\left((1-\lambda_{k+1})\sum_{i=1}^{k}\frac{\lambda_i}{1-\lambda_{k+1}}x_i+\lambda_{k+1}x_{k+1}\right)$$
$$\leq(1-\lambda_{k+1})\sum_{i=1}^{k}\frac{\lambda_i}{1-\lambda_{k+1}}f(x_i)+\lambda_{k+1}f(x_{k+1})=\sum_{i=1}^{k+1}\lambda_i f(x_i).$$
这说明不等式对$n=k+1$也成立,从而定理得证.

\begin{ti}
\textbf{定理4.16(L'Hospital法则)} \\
(1)(之一)设$f,g$在$(a,b)$中可导,且$g'$处处非零.如果$\displaystyle\lim_{x\to a^+}f(x)=\lim_{x\to a^+}g(x)=0$,且$\displaystyle\lim_{x\to a^+}\dfrac{f'(x)}{g'(x)}=\alpha$,其中$\alpha$为实数或$\pm\infty$,则$\displaystyle\lim_{x\to a^+}\dfrac{f(x)}{g(x)}=\alpha$.\\
(2)(之二)设$f,g$在$(a,b)$中可导,且$g'$处处非零.如果$\displaystyle\lim_{x\to a^+}g(x)=\infty$,且$\displaystyle\lim_{x\to a^+}\dfrac{f'(x)}{g'(x)}=\alpha$,其中$\alpha$为实数或$\pm\infty$,则$\displaystyle\lim_{x\to a^+}\dfrac{f(x)}{g(x)}=\alpha$.
\end{ti}

\textbf{证明}:(1)补充定义$f(a)=g(a)=0$,则$f,g\in C^0[a,b)$.由$Cauchy$中值定理,任给$x\in(a,b)$,存在$\xi\in(a,x)$,使得
$$\frac{f(x)}{g(x)}=\frac{f(x)-f(a)}{g(x)-g(a)}=\frac{f'(\xi)}{g'(\xi)}.$$
当$x\to a^+$时,$\xi\to a^+$,从而$\displaystyle\lim_{x\to a^+}\frac{f(x)}{g(x)}=\lim_{x\to a^+}\frac{f'(\xi)}{g'(\xi)}=\alpha$. \\
(2)分情况讨论.\\
\romannum{1}  当$\alpha=0$时.
根据题设,任给$\varepsilon>0$,存在$\eta>0$,当$x\in(a,a+\eta)$时
$$\left\lvert\dfrac{f'(x)}{g'(x)}\right\rvert<\dfrac{\varepsilon}{2}.$$
此时,根据$Cauchy$中值定理,存在$\xi\in(x,a+\eta)$,使得
$$\dfrac{f(x)-f(a+\eta)}{g(x)-g(a+\eta)}=\dfrac{f'(\xi)}{g'(\xi)}$$.上式可以改写为
$$\frac{f(x)}{g(x)}=\frac{f'(\xi)}{g'(\xi)}\left[1-\frac{g(a+\eta)}{g(x)}\right]+\frac{f(a+\eta)}{g(x)}.$$
利用$\left\lvert\dfrac{f'(x)}{g'(x)}\right\rvert<\dfrac{\varepsilon}{2}$以及$\displaystyle\lim_{x\to a^+}g(x)=\infty$不难得知,
存在正数$\delta<\eta$,使得当$x\in(a,a+\delta)$时
$$\left\lvert\dfrac{f(x)}{g(x)}\right\rvert<\varepsilon,$$这说明$\displaystyle\lim_{x\to a^+}\dfrac{f(x)}{g(x)}=0$.\\
\romannum{2}  当$\alpha\in\mathbb{R}$时,通过将$f$换成$f-\alpha g$可以转化为$\alpha=0$的情形;\\
\romannum{3}  当$\alpha=\pm\infty$时与$\alpha=0$情形的证明类似.

\begin{ti}
\textbf{定理4.17(Taylor公式,带Peano余项)}  设函数$f$在$x_0$处有$n$阶导数,则当$x\to x_0$时
$$f(x)=\sum_{k=0}^{n}\frac{f^{(k)}(x_0)}{k!}(x-x_0)^{k}+o\bigl((x-x_0)^{n}\bigr),$$
其中我们约定$f^{(0)}=f$,$(x-x_0)^{0}=1$.
\end{ti}

\textbf{证明}:记$R_n(x)=f(x)-\displaystyle\sum_{k=0}^{n}\frac{f^{(k)}(x_0)}{k!}(x-x_0)^k$.\\
由$f$在$x_0$处$n$阶可导,对$\dfrac{R_n(x)}{(x-x_0)^n}$反复应用$L'Hospital$法则($n-1$次),可得
$$\lim_{x\to x_0}\frac{R_n(x)}{(x-x_0)^n}=\lim_{x\to x_0}\frac{R_n^{(n-1)}(x)}{n!(x-x_0)}=\frac{1}{n!}\lim_{x\to x_0}\left[\frac{f^{(n-1)}(x)-f^{(n-1)}(x_0)}{x-x_0}-f^{(n)}(x_0)\right]=0,$$
其中最后一个等号用到$f^{(n)}(x_0)$的存在性.

\begin{ti}
\textbf{定理4.18(Taylor公式,带Lagrange余项)}  设$f$为$n$阶可导函数,则
$$f(x)=\sum_{k=0}^{n-1}\frac{f^{(k)}(x_0)}{k!}(x-x_0)^{k}+\frac{f^{(n)}(\xi)}{n!}(x-x_0)^{n},$$
其中$\xi=x_0+\theta(x-x_0)$,$\theta\in(0,1)$.
\end{ti}

\textbf{证明}:考虑以$t$为变量的函数$F(t)=\displaystyle\sum_{k=0}^{n-1}\frac{f^{(k)}(t)}{k!}(x-t)^k$.对$t$求导可得
$$F'(t)=\sum_{k=0}^{n-1}\left[\frac{f^{(k+1)}(t)}{k!}(x-t)^k-\frac{f^{(k)}(t)}{(k-1)!}(x-t)^{k-1}\right]=\frac{f^{(n)}(t)}{(n-1)!}(x-t)^{n-1}.$$
根据$F$的构造,我们有$F(x)-F(x_0)=f(x)-\displaystyle\sum_{k=0}^{n-1}\frac{f^{(k)}(x_0)}{k!}(x-x_0)^k$.\\
由$Lagrange$微分中值定理,存在$\xi=x_0+\theta(x-x_0)\,(\theta\in(0,1))$,使得
$$f(x)-\sum_{k=0}^{n-1}\frac{f^{(k)}(x_0)}{k!}(x-x_0)^k=F'(\xi)(x-x_0)=\frac{f^{(n)}(\xi)}{n!}(x-x_0)^n.$$

\begin{ti}
\textbf{定义4.19(Riemann可积)}  设$f$为$[a,b]$中的有界函数.如果存在实数$I$,使得任给$\varepsilon>0$,均存在$\delta>0$,当分割$P$的模满足$\left\lVert P\right\rVert<\delta$时,均有
$$\left\lvert\sum_{i=1}^{n}f(\xi_i)\Delta x_i-I\right\rvert<\varepsilon,\quad \forall\,\xi_i\in[x_{i-1},x_i],$$
则称$f$在$[a,b]$中\textbf{Riemann可积}(简称可积),记为$f\in R[a,b]$. $I$称为$f$在$[a,b]$中的\textbf{Riemann积分}(简称积分),记为
$$\int_a^b f(x)\,\mathrm{d}x=I=\lim_{\left\lVert P\right\rVert\to 0}\sum_{i=1}^{n}f(\xi_i)\Delta x_i.$$
\end{ti}

\begin{ti}
\textbf{定理4.20(可积的充要条件)}  设函数$f$在$[a,b]$上有界,则以下命题等价: \\
(1) $f$在$[a,b]$中$Riemann$可积; \\
(2) $f$在$[a,b]$中上积分和下积分相等; \\
(3) $\displaystyle\lim_{\left\lVert P\right\rVert\to 0}\sum_{i=1}^{n}\omega_i\varDelta x_i=0$; \\
(4) 任给$\varepsilon>0,$存在$[a,b]$中某个分割$P$,使得
$$S(P)-s(P)=\sum_{i=1}^{n}\omega_i\varDelta x_i<\varepsilon$$
\end{ti}

\begin{ti}
\textbf{定义4.21(零测集)}  设$A\subset\mathbb{R}$.如果任给$\varepsilon>0$,均可找到至多可数个开区间$\{I_i\}$,使得$A$包含于这些区间之并,且
$$\sum_{i\leq n}\left\lvert I_i\right\rvert\leq\varepsilon,\quad \forall\,n\geq 1,$$
即这些区间的长度之和不超过$\varepsilon$,则称$A$为\textbf{零测集}.
\end{ti}

\begin{ti}
\textbf{定理4.22(Lebesgue可积性判别法)}  设函数$f$在$[a,b]$上有界,则$f$在$[a,b]$上Riemann可积当且仅当$f$的间断点集为零测集.
\end{ti}

\begin{ti}
\textbf{命题4.23}  设$f\in C^0[a,b]$,$c\in[a,b]$.定义
$$g(x)=\int_c^x f(t)\,\mathrm{d}t,\quad x\in[a,b],$$
则$g'=f$在$[a,b]$中处处成立.
\end{ti}

\begin{ti}
\textbf{定理4.24(积分第二中值定理)}  设函数$f$在$[a,b]$上可积,$g$在$[a,b]$上单调,则存在$\xi\in[a,b]$,使得
$$\int_a^b f(x)g(x)\,\mathrm{d}x=g(a)\int_a^{\xi}f(x)\,\mathrm{d}x+g(b)\int_{\xi}^{b}f(x)\,\mathrm{d}x$$.
\end{ti}

\textbf{证明}:记$F(x)=\displaystyle\int_a^x f(t)\,\mathrm{d}t$,则$F$在$[a,b]$上连续.先设$g$单调递减且$g(b)=0$(否则考虑$g-g(b)$).任取$[a,b]$的分割$P:a=x_0<x_1<\cdots<x_n=b$,利用$g$的单调性、$F$的连续性以及$Abel$求和,可以逐步选取中值点,使得存在$\xi_P\in[a,b]$,满足
$$\int_a^b f(x)g(x)\,\mathrm{d}x=F(b)g(b)-F(a)g(a)-F(\xi_P)[g(b)-g(a)].$$
取一列分割$\{P_n\}$,使$\left\lVert P_n\right\rVert\to 0$,得数列$\{\xi_n\}\subset[a,b]$.根据$Bolzano$定理,$\{\xi_n\}$存在收敛子列,不妨设它自身收敛于$\xi\in[a,b]$,则
$$\int_a^b f(x)g(x)\,\mathrm{d}x=F(b)g(b)-F(a)g(a)-F(\xi)[g(b)-g(a)]=g(a)[F(\xi)-F(a)]+g(b)[F(b)-F(\xi)]$$
$\displaystyle=g(a)\int_a^{\xi}f(x)\,\mathrm{d}x+g(b)\int_{\xi}^{b}f(x)\,\mathrm{d}x.$\\

证毕.

\section{多元微分学与多元积分学}
\begin{ti}
\textbf{定义5.1(方向导数)}  设函数$f$在点$\mathbf{x}_0\in\mathbb{R}^n$的某邻域中有定义,$\mathbf{v}$为单位向量.如果极限
$$\lim_{t\to 0}\frac{f(\mathbf{x}_0+t\mathbf{v})-f(\mathbf{x}_0)}{t}$$
存在,则称其为$f$在$\mathbf{x}_0$处沿方向$\mathbf{v}$的\textbf{方向导数},记为$\dfrac{\partial f}{\partial\mathbf{v}}(\mathbf{x}_0)$.
\end{ti}

\begin{ti}
\textbf{定义5.2(微分)}  设$D\subset \mathbb{R}^n$为开集,$f:D\rightarrow \mathbb{R}$为多元函数,
$x^0\in D$.如果存在线性映射$L:\mathbb{R}^n\rightarrow\mathbb{R}$,使得在$x^0$附近成立
$$f(x)-f(x^0)=L(x-x^0)+o(\left\lVert x-x^0\right\rVert)\;(x\rightarrow x^0),$$
成立,则称$f$在$x^0$处\textbf{可微},称$L$为$f$在$x^0$处的\textbf{微分},也记为$df(x^0)$
\end{ti}

\begin{ti}
\textbf{命题5.3}  设函数$f$在$\mathbf{x}_0$处可微,则$f$在$\mathbf{x}_0$处连续,且$f$在$\mathbf{x}_0$处的方向导数都存在,并且有
$$\frac{\partial f}{\partial \mathbf{u}}(\mathbf{x}_0)=df(\mathbf{x}_0)(\mathbf{u}),$$
其中$\mathbf{u}$为单位向量.
\end{ti}

\begin{ti}
\textbf{定理5.4(可微的充分条件)}  设函数$f$在$\mathbf{x}_0$的某邻域内各偏导数均存在且在$\mathbf{x}_0$处连续,则$f$在$\mathbf{x}_0$处可微.
\end{ti}

\begin{ti}
\textbf{定义5.5(梯度)}  设函数$f$在$\mathbf{x}_0$处可微,记
$$\nabla f(\mathbf{x}_0)=\left(\frac{\partial f}{\partial x_1}(\mathbf{x}_0),\ldots,\frac{\partial f}{\partial x_n}(\mathbf{x}_0)\right),$$
称为$f$在$\mathbf{x}_0$处的\textbf{梯度}.利用$\mathbb{R}^n$中的标准内积,有$df(\mathbf{x}_0)(\mathbf{v})=\nabla f(\mathbf{x}_0)\cdot\mathbf{v}$.
\end{ti}

\begin{ti}
\textbf{命题5.6(切平面)}  设曲面$S$由方程$F(x,y,z)=0$给出,$F$可微且$\nabla F(P_0)\neq\mathbf{0}$,则$S$在点$P_0=(x_0,y_0,z_0)$处的切平面方程为
$$F_x(P_0)(x-x_0)+F_y(P_0)(y-y_0)+F_z(P_0)(z-z_0)=0$$.
\end{ti}

\begin{ti}
\textbf{定理5.7(链式法则)}  设$D$和$\varDelta$分别为$\mathbb{R}^n\text{和}\mathbb{R}^m$
中的开集,$f:D\rightarrow \mathbb{R}^m\text{和}g:\varDelta\rightarrow\mathbb{R}^l$为向量值函数,且$f(D)\subset\varDelta.$
如果$f$在$x^0\in D$处可微,$g$在$y^0=f(x^0)$处可微,则复合函数$h=g\circ f$在$x^0$处可微,且
$$Jh(x^0)=Jg(y^0)\cdot Jf(x^0)$$.
\end{ti}

\textbf{证明}:由$f$在$x^0$处可微可知,在$x^0$附近成立$f(x)-f(x^0)=Jf(x^0)(x-x^0)+o(\left\lVert x-x^0\right\rVert)$.这说明存在常数$C$,使得$\left\lVert f(x)-f(x^0)\right\rVert\leq C\left\lVert x-x^0\right\rVert$.特别地,当$x\to x^0$时,$f(x)\to f(x^0)$.同理,$g$在$y^0$处可微,故$g(y)-g(y^0)=Jg(y^0)(y-y^0)+o(\left\lVert y-y^0\right\rVert)$.在上式中代入$y=f(x)$可得
$$h(x)-h(x^0)=Jg(y^0)[f(x)-f(x^0)]+o(\left\lVert f(x)-f(x^0)\right\rVert)$$
$$=[Jg(y^0)\cdot Jf(x^0)](x-x^0)+o(\left\lVert x-x^0\right\rVert),$$
这说明$h$在$x^0$处可微,且$Jh(x^0)=Jg(y^0)\cdot Jf(x^0)$.

\begin{ti}
\textbf{定理5.8(微分中值定理)} 设$D\subset \mathbb{R}^n$为凸域,$f:D\rightarrow\mathbb{R}$在$D$中处处可微,则任给$x,y\in D$,
存在$\theta\in(0,1)$,使得
$$f(x)-f(y)=\nabla f(\xi)\cdot(x-y),\;\xi=\theta x+(1-\theta)y$$.
\end{ti}

\textbf{证明}:令$\sigma(t)=tx+(1-t)y$,由$D$为凸域可知当$t\in[0,1]$时$\sigma(t)\in D$.对一元函数$\varphi(t)=f(\sigma(t))$用微分中值定理可知存在$\theta\in(0,1)$,使得$\varphi(1)-\varphi(0)=\varphi'(\theta)$.由链式法则可得
$$f(x)-f(y)=\varphi(1)-\varphi(0)=\varphi'(\theta)=\nabla f(\xi)\cdot(x-y),$$
其中$\xi=\sigma(\theta)=\theta x+(1-\theta)y$.由$f(x)=\varphi(1)$,$f(y)=\varphi(0)$可知欲证结论成立.

\begin{ti}
\textbf{定理5.9(拟微分中值定理)}  设$D\subset\mathbb{R}^n$为凸域,$f:D\rightarrow\mathbb{R}^m$在$D$中处处可微,则对任意的$x,y\in D$,存在$\xi\in D$,使得
$$\left\lVert f(x)-f(y)\right\rVert\leq\left\lVert Jf(\xi)\right\rVert\cdot\left\lVert x-y\right\rVert.$$
\end{ti}

\textbf{证明}:基本的想法是对$f$的分量的线性组合应用微分中值定理.为此,不妨设$f(x)\neq f(y)$.
任意取定$\mathbb{R}^m$中的单位向量$u=(u_1,\ldots,u_m)$,记
$$g=u\cdot f=\displaystyle\sum_{i=1}^{m}u_if_i,$$
则$g$为$D$中可微函数.根据微分中值定理,存在$\xi\in D$,使得
$$g(x)-g(y)=\nabla g(\xi)\cdot(x-y).$$
注意到$\nabla g(\xi)=\displaystyle\sum_{i=1}^{m}u_i\nabla f_i(\xi)$,利用$Cauchy-Schwarz$不等式可得
$$\left\lVert\nabla g(\xi)\right\rVert\leq\left\lVert Jf(\xi)\right\rVert.$$
由$g(x)-g(y)=u\cdot[f(x)-f(y)]$可得
$$\left\lvert u\cdot[f(x)-f(y)]\right\rvert\leq\left\lVert Jf(\xi)\right\rVert\cdot\left\lVert x-y\right\rVert.$$
在上式中取$u=\dfrac{f(x)-f(y)}{\left\lVert f(x)-f(y)\right\rVert}$就完成了定理的证明.

%------------------------------------
%上次检查到这里，应该从定理5.10开始检查
%---------------------------------------
\begin{ti}
\textbf{定理5.10(逆映射定理)}  设$D\subset\mathbb{R}^n$为开集,$f:D\rightarrow\mathbb{R}^n$为$C^k\,(k\geq 1)$映射,$\mathbf{x}_0\in D$.如果$\det Jf(\mathbf{x}_0)\neq 0$,则存在$\mathbf{x}_0$的开邻域$U\subset D$以及$y_0=f(\mathbf{x}_0)$的开邻域$V\subset\mathbb{R}^n$,使得$f|_U:U\rightarrow V$是可逆映射,且其逆仍为$C^k$映射.
\end{ti}

\textbf{证明}:不失一般性,可设$\mathbf{x}_0=0$,$y_0=0$.以$L$记$f$在$\mathbf{x}_0=0$处的微分,则$L$可逆,且$L^{-1}\circ f$在$\mathbf{x}_0$处的微分为恒同映射.如果欲证结论对$L^{-1}f$成立,则对$f$也成立.因此,不妨从一开始就假设$Jf(\mathbf{x}_0)=I_n$.在$\mathbf{x}_0=0$附近,$f$是恒同映射的小扰动:$f(x)=x+g(x)$,$Jg(0)=0$,扰动项$g(x)=f(x)-x$为$C^k$映射.因此,存在$\delta>0$使得$\left\lVert Jg(x)\right\rVert\leq\dfrac{1}{2}$对一切$x\in B_\delta(0)\subset D$成立.由拟微分中值定理可知$\left\lVert g(x_1)-g(x_2)\right\rVert\leq\dfrac{1}{2}\left\lVert x_1-x_2\right\rVert$,$\forall x_1,x_2\in B_\delta(0)$.给定$y\in B_{\delta/2}(0)$,我们在$B_\delta(0)$中解方程$f(x)=y$,即$x=y-g(x)$.记$\varphi(x)=y-g(x)$,则$\varphi$为$B_\delta(0)$到自身的压缩映射.根据压缩映射原理,方程在$B_\delta(0)$中有唯一解,记为$x_y$.令$U=f^{-1}(B_{\delta/2}(0))\cap B_\delta(0)$,$V=B_{\delta/2}(0)$,则$f|_U:U\to V$是一一映射,其逆映射$h(y)=x_y$满足$y-g(h(y))=h(y)$. \\
(1)$h:V\to U$是连续映射:当$y_1,y_2\in V$时,
$$\left\lVert h(y_1)-h(y_2)\right\rVert\leq\left\lVert y_1-y_2\right\rVert+\frac{1}{2}\left\lVert h(y_1)-h(y_2)\right\rVert,$$
这说明$\left\lVert h(y_1)-h(y_2)\right\rVert\leq 2\left\lVert y_1-y_2\right\rVert$. \\
(2)$h:V\to U$是可微映射:设$y_0\in V$,则对$y\in V$,有
$$h(y)-h(y_0)=(y-y_0)-[g(h(y))-g(h(y_0))]=(y-y_0)-Jg(h(y_0))\cdot(h(y)-h(y_0))+o(\left\lVert h(y)-h(y_0)\right\rVert).$$
利用$Jf=I_n+Jg$和(1),上式可改写为$Jf(h(y_0))\cdot(h(y)-h(y_0))=(y-y_0)+o(\left\lVert y-y_0\right\rVert)$,因而$h(y)-h(y_0)=[Jf(h(y_0))]^{-1}\cdot(y-y_0)+o(\left\lVert y-y_0\right\rVert)$.这说明$h$在$y_0$处可微. \\
(3)$h:V\to U$为$C^k$映射:由(2)可知$Jh(y)=[Jf(h(y))]^{-1}$,$\forall y\in V$.由$f\in C^k$知$Jf\in C^{k-1}$.由(2)及上式可推出$Jh$连续,即$h\in C^1$.再由$Jf\in C^{k-1}$,$h\in C^1$及上式可推出$Jh\in C^1$,即$h\in C^2$.依此类推,最后我们就得到$h\in C^k$.

\begin{ti}
\textbf{定理5.11(隐映射定理)}  设$W\subset \mathbb{R}^{m+n}$为开集,$W$中的点用$(x,y)$表示,
其中$x=(x_1,x_2,...,x_n),y=(y_1,y_2,...,y_m).f:W\rightarrow\mathbb{R}^m\text{为}C^k$映射,用分量表示为
$$f(x,y)=(f_1(x,y),f_2(x,y),...,f_m(x,y)).$$
设$(x^0,y^0)\in W,\text{且}det J_yf(x^0,y^0)\neq 0,\text{其中}J_yf(x,y)=\left(\dfrac{\partial f_i}{\partial y_j}(x,y)\right)_{m\times m}$.
则存在$x^0$的开邻域$V\subset\mathbb{R}^n$以及唯一的$C^k$映射$\psi:V\rightarrow\mathbb{R}^m,$使得 \\
(1)$y^0=\psi(x^0),f(x,\psi(x))=f(x^0,y^0),\forall\,x\in V$; \\
(2)$J\psi(x)=-[J_yf(x,\psi(x))]^{-1}J_xf(x,\psi(x)),\text{其中}J_xf(x,y)=\left(\dfrac{\partial f_i}{\partial x_j}(x,y)\right)_{m\times n}$.
\end{ti}

\textbf{证明}:定义$F:W\to\mathbb{R}^{n+m}$为$F(x,y)=(x,f(x,y))$.在$(x^0,y^0)$处,
$$JF(x^0,y^0)=\begin{pmatrix}I_n&0\\J_xf(x^0,y^0)&J_yf(x^0,y^0)\end{pmatrix},$$
故$\det JF(x^0,y^0)=\det J_yf(x^0,y^0)\neq 0$.对$F$在$(x^0,y^0)$处利用逆映射定理,存在$(x^0,y^0)$的开邻域$\widetilde W\subset W$以及$F$的$C^k$逆映射,使得$F^{-1}(x,f(x^0,y^0))=(x,\psi(x))$在$x^0$的某开邻域$V$上有定义,且$\psi$为$C^k$映射.由$F$的定义可得$(x,f(x,\psi(x)))=(x,f(x^0,y^0))$,即$y^0=\psi(x^0)$,$f(x,\psi(x))=f(x^0,y^0)$对一切$x\in V$成立.对$f(x,\psi(x))=f(x^0,y^0)$关于$x$求导,由链式法则可得$J_xf(x,\psi(x))+J_yf(x,\psi(x))\cdot J\psi(x)=0$,即$J\psi(x)=-[J_yf(x,\psi(x))]^{-1}J_xf(x,\psi(x))$.唯一性由$F$的可逆性给出.

\begin{ti}
\textbf{定义5.12(极值)}  设函数$f$定义在$\mathbf{x}_0$的某邻域内.如果存在$\delta>0$,使得当$\left\lVert\mathbf{x}-\mathbf{x}_0\right\rVert<\delta$时$f(\mathbf{x})\leq f(\mathbf{x}_0)$恒成立,
则称$\mathbf{x}_0$为$f$的\textbf{极大值点},$f(\mathbf{x}_0)$为\textbf{极大值};极小值点与极小值类似定义.
\end{ti}

\begin{ti}
\textbf{定理5.13(极值的必要条件)}  设$\mathbf{x}_0$为函数$f$的极值点.如果$x_0$为$D$的内点,且$f$在$\mathbf{x}_0$处可微,则$\nabla f(\mathbf{x}_0)=\mathbf{0}$.
\end{ti}

\textbf{证明}:任取单位向量$u\in\mathbb{R}^n$,考虑一元函数$\varphi(t)=f(\mathbf{x}_0+tu)$.因为$\mathbf{x}_0$为内点,当$|t|$充分小时,$\mathbf{x}_0+tu\in D$.于是$t=0$为$\varphi$的极值点.根据一元函数的$Fermat$定理和链式法则可得
$$0=\varphi'(0)=\nabla f(\mathbf{x}_0)\cdot u.$$
由$u$的任意性可知$\nabla f(\mathbf{x}_0)=\mathbf{0}$.

\begin{ti}
\textbf{定理5.14(极值的充分条件)}  设函数$f$在$\mathbf{x}_0$的某邻域内二阶连续可微,且$\nabla f(\mathbf{x}_0)=\mathbf{0}$,
记$Hesse$矩阵$\nabla^2 f=\left(\dfrac{\partial^2 f}{\partial x_i\partial x_j}\right)_{n\times n}$.\\
(1)如果$\nabla^2 f$正定,则$\mathbf{x}_0$为严格极小值点;\\
(2)如果$\nabla^2 f$负定,则$\mathbf{x}_0$为严格极大值点;\\
(3)如果$\nabla^2 f$不定,则$\mathbf{x}_0$不是极值点.
\end{ti}

\textbf{证明}:设$u=x-\mathbf{x}_0$,对一元函数$\varphi(t)=f(\mathbf{x}_0+tu)$应用$Taylor$公式可得
$$f(x)=\varphi(1)=\varphi(0)+\varphi'(0)+\frac{1}{2}\varphi''(\theta)=f(\mathbf{x}_0)+\frac{1}{2}(x-\mathbf{x}_0)\nabla^2 f(\xi)(x-\mathbf{x}_0)^T,$$
其中$\xi=\mathbf{x}_0+\theta u$,$\theta\in(0,1)$.由$\nabla^2 f$的连续性,当$x$充分接近$\mathbf{x}_0$时,$\nabla^2 f(\xi)$与$\nabla^2 f(\mathbf{x}_0)$保持同号性.因此,若$\nabla^2 f(\mathbf{x}_0)$正定,则对充分接近$\mathbf{x}_0$的$x\neq\mathbf{x}_0$有$f(x)>f(\mathbf{x}_0)$,即$\mathbf{x}_0$为严格极小值点;若$\nabla^2 f(\mathbf{x}_0)$负定,则$\mathbf{x}_0$为严格极大值点;若$\nabla^2 f(\mathbf{x}_0)$不定,则可选取方向$u$,使$\dfrac{1}{2}(x-\mathbf{x}_0)\nabla^2 f(\xi)(x-\mathbf{x}_0)^T$可正可负,故$\mathbf{x}_0$不是极值点.

\begin{ti}
\textbf{定理5.15(Lagrange乘数法)}  设$f$与约束函数$\Phi=(g_1,\ldots,g_m)$在$\mathbf{x}_0$附近为$C^1$映射,$\mathbf{x}_0\in\Sigma=\left\{\mathbf{x}\in\mathbb{R}^n\;\middle|\;\Phi(\mathbf{x})=0\right\}$为$f$的条件极值点.如果$J\Phi(\mathbf{x}_0)$的秩为$m$,则存在$\lambda^0=(\lambda_1,\ldots,\lambda_m)\in\mathbb{R}^m$,使得
$$\nabla f(\mathbf{x}_0)=\sum_{j=1}^{m}\lambda_j\nabla g_j(\mathbf{x}_0)$$.
\end{ti}

\textbf{证明}:不妨设$\dfrac{\partial(g_1,g_2,\ldots,g_m)}{\partial(x_1,x_2,\ldots,x_m)}(\mathbf{x}_0)\neq 0$.由隐映射定理,存在$\mathbf{z}_0=(x_{m+1}^0,\ldots,x_n^0)$的开邻域$V$以及$C^1$映射$\psi:V\to\mathbb{R}^m$,使得$\mathbf{y}_0=\psi(\mathbf{z}_0)$,$\Phi(\psi(\mathbf{z}),\mathbf{z})=0$对一切$\mathbf{z}\in V$成立,其中$\mathbf{y}=(x_1,\ldots,x_m)$,$\mathbf{z}=(x_{m+1},\ldots,x_n)$,$\mathbf{x}=(\mathbf{y},\mathbf{z})$.在$\mathbf{x}_0=(\mathbf{y}_0,\mathbf{z}_0)$处对$\Phi(\psi(\mathbf{z}),\mathbf{z})=0$求导,有$J_y\Phi(\mathbf{x}_0)\cdot J\psi(\mathbf{z}_0)+J_z\Phi(\mathbf{x}_0)=0$,即$J\psi(\mathbf{z}_0)=-(J_y\Phi)^{-1}(\mathbf{x}_0)\cdot J_z\Phi(\mathbf{x}_0)$.因为$\mathbf{x}_0$为$f$的条件极值点,$\mathbf{z}_0$为$F(\mathbf{z})=f(\psi(\mathbf{z}),\mathbf{z})$的极值点,由极值的必要条件,$\nabla F(\mathbf{z}_0)=0$,即$\nabla f(\mathbf{x}_0)\cdot J\psi(\mathbf{z}_0)$与$\nabla f(\mathbf{x}_0)$在$\mathbf{z}$方向的分量之和为零.整理可得$\nabla f(\mathbf{x}_0)-\lambda^0\cdot J\Phi(\mathbf{x}_0)=0$对某个$\lambda^0=(\lambda_1,\ldots,\lambda_m)\in\mathbb{R}^m$成立,即$\nabla f(\mathbf{x}_0)=\displaystyle\sum_{j=1}^{m}\lambda_j\nabla g_j(\mathbf{x}_0)$.

\begin{ti}
\textbf{定义5.16(二重Riemann积分)}  设函数$f$在有界闭区域$D\subset\mathbb{R}^2$上有界.对$D$的任意分割$P=\{D_1,\ldots,D_n\}$,记
$$S(P)=\sum_{i=1}^{n}M_i\,\sigma(D_i),\qquad s(P)=\sum_{i=1}^{n}m_i\,\sigma(D_i)$$
其中$\displaystyle M_i=\sup_{D_i}f$,$\displaystyle m_i=\inf_{D_i}f$, $\sigma(D_i)$为$D_i$的面积.
如果$\displaystyle\inf_P S(P)=\sup_P s(P)$,则称$f$在$D$上\textbf{Riemann可积},
该公共值记为$\displaystyle\iint_D f(x,y)\,\mathrm{d}x\mathrm{d}y$.
\end{ti}

\begin{ti}
\textbf{命题5.17(重积分的交换积分次序)} 设$A\subset\mathbb{R}^2$为可求面积的有界集合,$f:A\rightarrow\mathbb{R}$为连续有界函数.
记$A$在$x$轴上的垂直投影为
$$I=\{x\in\mathbb{R}\;|\;\text{存在}y\text{使得}(x,y)\in A\}$$.
如果对于每一点$x\in I,A_x=\{y\in\mathbb{R}\;|\;(x,y)\in A\}$是区间(可退化为一点),则
$$\int_A f=\int_I dx\int_{A_x}f(x,y)dy.$$
\end{ti}

\begin{ti}
\textbf{命题5.18(重积分变量替换公式)}  设映射$\Phi:(u,v)\mapsto(x,y)$为有界闭区域$D'\to D$的$C^1$一一映射,且$\det D\Phi\neq 0$,函数$f$在$D$上连续,则
$$\iint_D f(x,y)\,\mathrm{d}x\mathrm{d}y=\iint_{D'}f(\Phi(u,v))\left\lvert\det D\Phi(u,v)\right\rvert\,\mathrm{d}u\mathrm{d}v.$$
特别地,极坐标变换$x=r\cos\theta, y=r\sin\theta$下,$\mathrm{d}x\mathrm{d}y=r\,\mathrm{d}r\mathrm{d}\theta$.
\end{ti}

\begin{ti}
\textbf{定义5.19(第一型曲线积分)}  设$\sigma:[\alpha,\beta]\rightarrow\mathbb{R}^n$为参数曲线,$f$是定义在$\sigma$上的有界函数.任给$[\alpha,\beta]$的分割$\pi:\alpha=t_0<t_1<\cdots<t_m=\beta$,取$\xi_i\in[t_{i-1},t_i]$,考虑和$\displaystyle\sum_{i=1}^{m}f(\sigma(\xi_i))\Delta s_i$,其中$\Delta s_i=s(t_i)-s(t_{i-1})$.如果极限
$$\lim_{\left\lVert\pi\right\rVert\to 0}\sum_{i=1}^{m}f(\sigma(\xi_i))\Delta s_i$$
存在且与$\{\xi_i\}$的选取无关,则称此极限为$f$在$\sigma$上的\textbf{第一型曲线积分},记为$\displaystyle\int_\sigma f\,\mathrm{d}s$.当$\sigma$为分段$C^1$曲线时,
$$\int_\sigma f\,\mathrm{d}s=\int_{\alpha}^{\beta}f(\sigma(t))\left\lVert\sigma'(t)\right\rVert\,\mathrm{d}t.$$
\end{ti}

\begin{ti}
\textbf{定义5.20(第二型曲线积分)} 对于分段$C^1$曲线,有
$$\int_{\sigma}f_1dx_1+f_2dx_2+...+f_ndx_n=\int_{\alpha}^{\beta}F(\sigma)\cdot\sigma'(t)\mathrm{d}t$$,
其中$F=(f_1,f_2,...,f_n),\sigma'(t)=(x_1'(t),x_2'(t),...,x_n'(t))$为曲线$\sigma$的切向量.
\end{ti}

\begin{ti}
\textbf{定义5.21(第一型曲面积分)}  设曲面$S$由参数方程$\mathbf{r}(u,v)$, $(u,v)\in D$给出,函数$f$在$S$上连续,则
$$\iint_S f\,\mathrm{d}S=\iint_D f(\mathbf{r}(u,v))\left\lVert\mathbf{r}_u\times\mathbf{r}_v\right\rVert\,\mathrm{d}u\mathrm{d}v$$.
\end{ti}

\begin{ti}
\textbf{定义5.22(第二型曲面积分)}  设向量场$\mathbf{F}=(P,Q,R)$在定向曲面$S$上连续,$\mathbf{n}$为$S$的单位法向量,则
$$\iint_S P\,\mathrm{d}y\wedge\mathrm{d}z+Q\,\mathrm{d}z\wedge\mathrm{d}x+R\,\mathrm{d}x\wedge\mathrm{d}y=\iint_S\mathbf{F}\cdot\mathbf{n}\,\mathrm{d}S$$.
\end{ti}

\begin{ti}
\textbf{定理5.23(Green公式)}  设$\Omega$为平面有界区域,其边界由有限条$C^1$曲线组成,边界的定向为诱导定向.如果$P,Q$为$\Omega$上的$C^1$函数,则
$$\oint_{\partial\Omega}P\,\mathrm{d}x+Q\,\mathrm{d}y=\iint_\Omega\left(\frac{\partial Q}{\partial x}-\frac{\partial P}{\partial y}\right)\mathrm{d}x\mathrm{d}y$$.
\end{ti}

\textbf{证明}:先考虑$\Omega$为矩形的特殊情形.设$\Omega=[a,b]\times[c,d]$,则由$Newton-Leibniz$公式,
$$\iint_\Omega\frac{\partial Q}{\partial x}\,\mathrm{d}x\mathrm{d}y=\int_c^d[Q(b,y)-Q(a,y)]\,\mathrm{d}y=\oint_{\partial\Omega}Q\,\mathrm{d}y,$$
$$\iint_\Omega\frac{\partial P}{\partial y}\,\mathrm{d}x\mathrm{d}y=\int_a^b[P(x,d)-P(x,c)]\,\mathrm{d}x=-\oint_{\partial\Omega}P\,\mathrm{d}x,$$
两式相减即得$Green$公式.对一般的有界区域,将其边界分为有限条$C^1$曲线段,将区域适当分割为有限个标准区域(或利用保定向的坐标变换化为矩形),在每个标准区域上应用已证结论并求和,交叉边界上的积分相互抵消,即得$Green$公式对$\Omega$成立.

\begin{ti}
\textbf{定理5.24(Gauss公式)}  设$\Omega$为$\mathbb{R}^3$中的有界区域,其边界由有限个$C^1$曲面组成,曲面的定向为诱导定向(即外侧方向).如果$P,Q,R$为$\Omega$上的连续可微函数,则
$$\iiint_\Omega\left(\frac{\partial P}{\partial x}+\frac{\partial Q}{\partial y}+\frac{\partial R}{\partial z}\right)\mathrm{d}x\mathrm{d}y\mathrm{d}z=\iint_{\partial\Omega}P\,\mathrm{d}y\wedge\mathrm{d}z+Q\,\mathrm{d}z\wedge\mathrm{d}x+R\,\mathrm{d}x\wedge\mathrm{d}y$$.
\end{ti}

\textbf{证明}:先考虑$\Omega$为长方体$[a,b]\times[c,d]\times[e,f]$的特殊情形.由$Newton-Leibniz$公式,
$$\iiint_\Omega\frac{\partial P}{\partial x}\,\mathrm{d}x\mathrm{d}y\mathrm{d}z=\iint_{[c,d]\times[e,f]}[P(b,y,z)-P(a,y,z)]\,\mathrm{d}y\mathrm{d}z=\iint_{\partial\Omega}P\,\mathrm{d}y\wedge\mathrm{d}z,$$
对$Q,R$同理,三式相加即得$Gauss$公式.对一般的有界区域,将其边界分为有限个$C^1$曲面片,将区域适当分割为有限个标准区域(或利用保定向的坐标变换),在每个标准区域上应用已证结论并求和,交叉边界上的积分相互抵消,即得$Gauss$公式对$\Omega$成立.

\begin{ti}
\textbf{定理5.25(Stokes公式)}  设$F=(P,Q,R)$,$S$为分片光滑的定向曲面,其边界$\partial S$为分段光滑的闭曲线(取诱导定向),向量场$\mathbf{F}$在含$S$的区域内$C^1$,则
$$\oint_{\partial S}\mathbf{F}\cdot T\mathrm{d}\mathbf{s}=\iint_S rot\,F\cdot\mathbf{n}\,\mathrm{d}\sigma.$$
其中$rot\,F=(\dfrac{\partial R}{\partial y}-\dfrac{\partial Q}{\partial z},\dfrac{\partial P}{\partial z}-\dfrac{\partial R}{\partial x},\dfrac{\partial Q}{\partial x}-\dfrac{\partial P}{\partial y})$
称为$F$的旋度,$T$为曲线的单位切向量,$s$为弧长参数.
\end{ti}

\textbf{证明}:设曲面$S$有$C^2$参数表示$\mathbf{r}(u,v)$,$(u,v)\in D$,区域$\Omega\subset D$与$S$对应,其边界$\partial\Omega$取诱导定向.由第二型曲线积分的定义,
$$\oint_{\partial S}\mathbf{F}\cdot T\,\mathrm{d}s=\oint_{\partial\Omega}\left[\mathbf{F}(\mathbf{r}(u,v))\cdot\mathbf{r}_u\right]\mathrm{d}u+\left[\mathbf{F}(\mathbf{r}(u,v))\cdot\mathbf{r}_v\right]\mathrm{d}v.$$
对上述平面曲线积分应用$Green$公式,得
$$\oint_{\partial S}\mathbf{F}\cdot T\,\mathrm{d}s=\iint_\Omega\left[\frac{\partial(\mathbf{F}\cdot\mathbf{r}_v)}{\partial u}-\frac{\partial(\mathbf{F}\cdot\mathbf{r}_u)}{\partial v}\right]\mathrm{d}u\mathrm{d}v=\iint_\Omega rot\,\mathbf{F}\cdot(\mathbf{r}_u\times\mathbf{r}_v)\,\mathrm{d}u\mathrm{d}v,$$
其中第二个等号由直接计算(用到$\mathbf{r}_{uv}=\mathbf{r}_{vu}$)得到.再由第一型曲面积分的定义,$\mathbf{n}\,\mathrm{d}\sigma=\mathbf{r}_u\times\mathbf{r}_v\,\mathrm{d}u\mathrm{d}v$,即得$\displaystyle\oint_{\partial S}\mathbf{F}\cdot T\,\mathrm{d}s=\iint_S rot\,\mathbf{F}\cdot\mathbf{n}\,\mathrm{d}\sigma$.

%------------------------
%声明:以下内容暂未经过审核
%-------------------------------
\section{级数}
\begin{ti}
\textbf{定义6.1(级数的收敛与发散)}  设$\{a_n\}$为数列,称形式和$\displaystyle\sum_{n=1}^{\infty}a_n=a_1+a_2+\cdots+a_n+\cdots$为\textbf{无穷级数},$a_n$称为通项或一般项,$S_n=\displaystyle\sum_{k=1}^{n}a_k$称为级数的第$n$个\textbf{部分和}.如果数列$\{S_n\}$收敛,则称级数$\displaystyle\sum_{n=1}^{\infty}a_n$\textbf{收敛},其和记为$\displaystyle\sum_{n=1}^{\infty}a_n=\lim_{n\to\infty}S_n$;否则就称级数$\displaystyle\sum_{n=1}^{\infty}a_n$\textbf{发散}.
\end{ti}

\begin{ti}
\textbf{定理6.2(收敛的必要条件)}  如果级数$\displaystyle\sum_{n=1}^{\infty}a_n$收敛,则$\displaystyle\lim_{n\to\infty}a_n=0$.
\end{ti}

\textbf{证明}:因为$a_n=S_n-S_{n-1}$,而部分和数列$\{S_n\}$收敛,故$\displaystyle\lim_{n\to\infty}a_n=\lim_{n\to\infty}(S_n-S_{n-1})=0$.

\begin{ti}
\textbf{定理6.3(Cauchy收敛准则)}  级数$\displaystyle\sum_{n=1}^{\infty}a_n$收敛当且仅当对任意的$\varepsilon>0$,存在$N$,使得当$m>n>N$时,
$$\left\lvert\sum_{k=n+1}^{m}a_k\right\rvert<\varepsilon$$.
\end{ti}

\textbf{证明}:级数$\displaystyle\sum a_n$收敛当且仅当部分和数列$\{S_n\}$收敛.由数列极限的$Cauchy$准则,$\{S_n\}$收敛当且仅当任给$\varepsilon>0$,存在$N$,当$m,n>N$时$|S_m-S_n|<\varepsilon$.而$S_m-S_n=a_{n+1}+a_{n+2}+\cdots+a_m$,即得欲证结论.

\begin{ti}
\textbf{定义6.4(绝对收敛与条件收敛)}  如果级数$\displaystyle\sum_{n=1}^{\infty}\left\lvert a_n\right\rvert$收敛,则称$\displaystyle\sum_{n=1}^{\infty}a_n$\textbf{绝对收敛};
如果$\displaystyle\sum_{n=1}^{\infty}a_n$收敛而$\displaystyle\sum_{n=1}^{\infty}\left\lvert a_n\right\rvert$发散,则称其\textbf{条件收敛}.
\end{ti}

\begin{ti}
\textbf{定理6.5}  如果级数$\displaystyle\sum_{n=1}^{\infty}a_n$绝对收敛,则它必收敛.
\end{ti}

\textbf{证明}:设级数$\displaystyle\sum a_n$绝对收敛,则任给$\varepsilon>0$,由$Cauchy$准则存在$N$,当$m>n>N$时$\displaystyle\sum_{k=n+1}^{m}|a_k|<\varepsilon$.于是$\left\lvert\displaystyle\sum_{k=n+1}^{m}a_k\right\rvert\leq\displaystyle\sum_{k=n+1}^{m}|a_k|<\varepsilon$,由$Cauchy$准则知级数$\displaystyle\sum a_n$收敛.

\begin{ti}
\textbf{定理6.6(比较判别法)}  设$0\leq a_n\leq b_n$对充分大的$n$恒成立.\\
(1)如果$\displaystyle\sum b_n$收敛,则$\displaystyle\sum a_n$也收敛;\\
(2)如果$\displaystyle\sum a_n$发散,则$\displaystyle\sum b_n$也发散.
\end{ti}

\textbf{证明}:设存在$M>0,N_0\geq 1$,使得当$n>N_0$时$a_n\leq Mb_n$.若$\displaystyle\sum b_n$收敛,则其部分和$\{B_n\}$有界,故$\{A_n\}$(从$N_0$之后)也有界,由基本判别法(正项级数收敛当且仅当部分和有界)知$\displaystyle\sum a_n$收敛;若$\displaystyle\sum a_n$发散,由上面已证的结论用反证法即得$\displaystyle\sum b_n$发散.

\begin{ti}
\textbf{定理6.7(比值判别法)}  设$a_n>0$,且$\displaystyle\lim_{n\to\infty}\frac{a_{n+1}}{a_n}=r$.\\
(1)当$r<1$时级数$\displaystyle\sum a_n$收敛;\\
(2)当$r>1$时级数$\displaystyle\sum a_n$发散.
\end{ti}

\textbf{证明}:若$r<1$,取$q$使$r<q<1$,则当$n$充分大时$\dfrac{a_{n+1}}{a_n}\leq q$,即$a_n\leq a_{N_0}q^{n-N_0}$,由比较判别法知$\displaystyle\sum a_n$收敛.若$r>1$,则当$n$充分大时$\dfrac{a_{n+1}}{a_n}\geq 1$,即$a_{n+1}\geq a_n$,通项不趋于$0$,级数发散.

\begin{ti}
\textbf{定理6.8(根值判别法)}  设$a_n\geq 0$,且$\displaystyle\lim_{n\to\infty}\sqrt[n]{a_n}=r$.\\
(1)当$r<1$时级数$\displaystyle\sum a_n$收敛;\\
(2)当$r>1$时级数$\displaystyle\sum a_n$发散.
\end{ti}

\textbf{证明}:若$r<1$,取$q$使$r<q<1$,则当$n$充分大时$\sqrt[n]{a_n}\leq q$,即$a_n\leq q^n$,由比较判别法(与几何级数$\displaystyle\sum q^n$比较)知$\displaystyle\sum a_n$收敛.若$r>1$,则对无穷多个$n$有$\sqrt[n]{a_n}\geq 1$,即$a_n\geq 1$,通项不趋于$0$,级数发散.

\begin{ti}
\textbf{定理6.9(积分判别法)}  设函数$f$在$[1,+\infty)$上非负单调递减,则级数$\displaystyle\sum_{n=1}^{\infty}f(n)$与广义积分$\displaystyle\int_1^{+\infty}f(x)\,\mathrm{d}x$同敛散.
\end{ti}

\textbf{证明}:由$f$单调递减可知,当$k\leq x\leq k+1$时,$f(k+1)\leq f(x)\leq f(k)$,故$f(k+1)\leq\displaystyle\int_k^{k+1}f(x)\,\mathrm{d}x\leq f(k)$.求和得
$$\sum_{k=1}^{n}f(k+1)\leq\int_1^{n+1}f(x)\,\mathrm{d}x\leq\sum_{k=1}^{n}f(k),$$
由此即知级数$\displaystyle\sum_{n=1}^{\infty}f(n)$与广义积分$\displaystyle\int_1^{+\infty}f(x)\,\mathrm{d}x$同敛散.

\begin{ti}
\textbf{定理6.10(Leibniz判别法)}  设$a_n$单调地趋于$0$,则级数$\displaystyle\sum_{n=1}^{\infty}(-1)^{n}a_n$收敛.
\end{ti}

\textbf{证明}:记$b_n=(-1)^n$,则级数$\displaystyle\sum b_n$的部分和显然有界.由$Dirichlet$判别法(设$\{a_n\}$单调地趋于$0$,级数$\displaystyle\sum b_n$的部分和有界,则级数$\displaystyle\sum a_nb_n$收敛)立得欲证结论.

\begin{ti}
\textbf{定义6.11(无穷乘积)}  设$\{p_n\}$为一列实数,称形式乘积$\displaystyle\prod_{n=1}^{\infty}p_n=p_1p_2\cdots p_n\cdots$为\textbf{无穷乘积},记$P_n=p_1p_2\cdots p_n$,称为\textbf{部分乘积}.如果数列$\{P_n\}$的极限存在,且极限为实数或正负无穷,则此极限称为无穷乘积的值,记为$\displaystyle\prod_{n=1}^{\infty}p_n=\lim_{n\to\infty}P_n$.当极限为非零实数时,称无穷乘积是\textbf{收敛}的,否则就称它是发散的.
\end{ti}

\begin{ti}
\textbf{定理6.12(无穷乘积的收敛判别法)}  设$a_n>-1$,则无穷乘积$\displaystyle\prod_{n=1}^{\infty}(1+a_n)$收敛当且仅当级数$\displaystyle\sum_{n=1}^{\infty}\ln(1+a_n)$收敛;
特别地,若$\displaystyle\sum_{n=1}^{\infty}\left\lvert a_n\right\rvert$收敛,则无穷乘积收敛.
\end{ti}

\textbf{证明}:设$p_n=1+a_n>0$.由$P_n=\displaystyle\prod_{k=1}^{n}p_k=\exp\left(\displaystyle\sum_{k=1}^{n}\ln p_k\right)$知,无穷乘积$\displaystyle\prod p_n$收敛当且仅当级数$\displaystyle\sum\ln(1+a_n)$收敛.特别地,若$\displaystyle\sum|a_n|$收敛,则当$n$充分大时$|\ln(1+a_n)|\leq 2|a_n|$(因为$\ln(1+x)\sim x$),由比较判别法知$\displaystyle\sum\ln(1+a_n)$绝对收敛,从而无穷乘积收敛.

\begin{ti}
\textbf{定理6.13(级数的乘积)}  设级数$\displaystyle\sum_{n=1}^{\infty}a_n$与$\displaystyle\sum_{n=1}^{\infty}b_n$绝对收敛,则它们的\textbf{Cauchy乘积}
$\displaystyle\sum_{n=1}^{\infty}c_n$(其中$c_n=\displaystyle\sum_{k=1}^{n}a_k b_{n+1-k}$)绝对收敛,且
$$\sum_{n=1}^{\infty}c_n=\left(\sum_{n=1}^{\infty}a_n\right)\left(\sum_{n=1}^{\infty}b_n\right)$$.
\end{ti}

\textbf{证明}:记$\displaystyle\sum|a_i|=A$,$\displaystyle\sum|b_j|=B$.对指标集$\{(i,j)\}$按任意次序求和,非负项级数$\displaystyle\sum_{i,j}|a_i||b_j|=\left(\sum|a_i|\right)\left(\sum|b_j|\right)=AB$收敛,故$\{a_ib_j\}$的任意重排(特别地,按$Cauchy$乘积的方式重排)绝对收敛.于是由重排定理,其和等于$\left(\displaystyle\sum a_n\right)\left(\displaystyle\sum b_n\right)$,即$Cauchy$乘积$\displaystyle\sum c_n$绝对收敛,且$\displaystyle\sum c_n=\left(\displaystyle\sum a_n\right)\left(\displaystyle\sum b_n\right)$.

\begin{ti}
\textbf{定理6.14(重排定理)}  绝对收敛级数的项可以任意重排而不改变其和.
\end{ti}

\textbf{证明}:设级数$\displaystyle\sum a_n$绝对收敛.记$a_n^+=\max\{a_n,0\}$,$a_n^-=\max\{-a_n,0\}$,则$\displaystyle\sum a_n^+$与$\displaystyle\sum a_n^-$均收敛.对原级数作任意重排$\displaystyle\sum a_{\sigma(n)}$,由非负项级数的重排不改变和,$\displaystyle\sum a_{\sigma(n)}^+=\sum a_n^+$,$\displaystyle\sum a_{\sigma(n)}^-=\sum a_n^-$,故$\displaystyle\sum a_{\sigma(n)}=\sum a_{\sigma(n)}^+-\sum a_{\sigma(n)}^-=\sum a_n^+-\sum a_n^-=\sum a_n$,且重排后的级数绝对收敛.

\begin{ti}
\textbf{定理6.15(Riemann重排定理)}  设级数$\displaystyle\sum a_n$条件收敛,则对任意的实数$S$(或$\pm\infty$),均可将其各项重排,使得重排后的级数收敛到$S$(或发散到$\pm\infty$).
\end{ti}

\textbf{证明}:设$\displaystyle\sum a_n$条件收敛.记$a_n^+=\max\{a_n,0\}$,$a_n^-=\max\{-a_n,0\}$,则$\displaystyle\sum a_n^+=\sum a_n^-=+\infty$(否则级数绝对收敛).设$\xi\in\mathbb{R}$.先依次取正项$a^+$中足够多的项,使部分和首次超过$\xi$(记所取项数为$m_1$);再依次取负项$a^-$中足够多的项,使部分和首次低于$\xi$(记所取项数为$n_1$);再取正项使部分和超过$\xi$,再取负项使部分和低于$\xi$,如此继续.由于$a_n\to 0$,每一步多取的项趋于$0$,故所得重排的部分和在$\xi$附近摆动且趋于$\xi$,即重排后的级数收敛于$\xi$.对发散到$\pm\infty$的情形可类似处理.

\begin{ti}
\textbf{定义6.16(一致收敛)}  设$\{f_n\}$为集合$E$上的函数列.如果对任意的$\varepsilon>0$,存在与$x$无关的$N$,使得当$n>N$时,
对一切$x\in E$均有$\left\lvert f_n(x)-f(x)\right\rvert<\varepsilon$,则称函数列$\{f_n\}$在$E$中\textbf{一致收敛}于$f$,记为$f_n\Rightarrow f$.
对于函数项级数$\displaystyle\sum f_n(x)$,若其部分和函数列一致收敛,则称该级数在$E$上\textbf{一致收敛}.
\end{ti}

\begin{ti}
\textbf{定理6.17(一致收敛的Cauchy准则)}  函数列$\{f_n\}$在$E$上一致收敛当且仅当对任意的$\varepsilon>0$,存在$N$,使得当$m,n>N$时,
对一切$x\in E$均有$\left\lvert f_m(x)-f_n(x)\right\rvert<\varepsilon$.
\end{ti}

\textbf{证明}:必要性显然.充分性:设函数列$\{f_n\}$满足条件,则对每个固定的$x_0\in E$,$\{f_n(x_0)\}$是$Cauchy$数列,从而是收敛数列,记其极限为$f(x_0)$,这样得到极限函数$f$.在条件$|f_m(x)-f_n(x)|<\varepsilon$中令$m\to\infty$,得$|f_n(x)-f(x)|\leq\varepsilon$对一切$n>N$和$x\in E$成立,即$\{f_n\}$一致收敛于$f$.

\begin{ti}
\textbf{定理6.18(Weierstrass判别法)}  设$\left\lvert f_n(x)\right\rvert\leq M_n$对一切$x\in E$成立,且正项级数$\displaystyle\sum M_n$收敛,
则函数项级数$\displaystyle\sum f_n(x)$在$E$上一致收敛.
\end{ti}

\textbf{证明}:由$\displaystyle\sum M_n$收敛,任给$\varepsilon>0$,存在$N$,当$m>n>N$时$\displaystyle\sum_{k=n+1}^{m}M_k<\varepsilon$.于是对一切$x\in E$,
$$\left\lvert\sum_{k=n+1}^{m}f_k(x)\right\rvert\leq\sum_{k=n+1}^{m}|f_k(x)|\leq\sum_{k=n+1}^{m}M_k<\varepsilon,$$
由一致收敛的$Cauchy$准则知$\displaystyle\sum f_n(x)$在$E$上一致收敛.

\begin{ti}
\textbf{定理6.19(一致收敛与连续性)}  设$\{f_n\}$在区间$I$中一致收敛于$f$.如果每一个$f_n$均为连续函数,则$f$也是连续函数.
\end{ti}

\textbf{证明}:任取$x_0\in I$,要证明$f$在$x_0$处连续.由题设,任给$\varepsilon>0$,存在$N=N(\varepsilon)$,当$n>N,x\in I$时,$|f_n(x)-f(x)|<\dfrac{\varepsilon}{3}$.取定$n_0=N+1$,由于$f_{n_0}$在$x_0$处连续,故存在$\delta=\delta(\varepsilon)>0$,使得当$|x-x_0|<\delta$时$|f_{n_0}(x)-f_{n_0}(x_0)|<\dfrac{\varepsilon}{3}$.于是当$|x-x_0|<\delta$时,
$$|f(x)-f(x_0)|\leq|f(x)-f_{n_0}(x)|+|f_{n_0}(x)-f_{n_0}(x_0)|+|f_{n_0}(x_0)-f(x_0)|<\frac{\varepsilon}{3}+\frac{\varepsilon}{3}+\frac{\varepsilon}{3}=\varepsilon,$$
这说明$f$在$x_0$处连续.由$x_0$的任意性,$f$在$I$上连续.

\begin{ti}
\textbf{定理6.20(逐项积分)}  设函数列$\{f_n\}$在$[a,b]$上连续且一致收敛到$f$,则
$$\lim_{n\to\infty}\int_a^b f_n(x)\,\mathrm{d}x=\int_a^b f(x)\,\mathrm{d}x$$.
\end{ti}

\textbf{证明}:因为$\{f_n\}$一致收敛于$f$,故$\displaystyle\sup_{x\in[a,b]}|f_n(x)-f(x)|\to 0$.于是
$$\left\lvert\int_a^b f_n(x)\,\mathrm{d}x-\int_a^b f(x)\,\mathrm{d}x\right\rvert\leq\int_a^b|f_n(x)-f(x)|\,\mathrm{d}x\leq(b-a)\sup_{x\in[a,b]}|f_n(x)-f(x)|\to 0,$$
即$\displaystyle\lim_{n\to\infty}\int_a^b f_n(x)\,\mathrm{d}x=\int_a^b f(x)\,\mathrm{d}x$.

\begin{ti}
\textbf{定理6.21(逐项求导)}  设函数列$\{f_n\}$在$[a,b]$上有连续导数,且$\{f_n\}$在$[a,b]$上收敛于$f$,$\{f'_n\}$在$[a,b]$上一致收敛,
则$f$在$[a,b]$上可导,且$\displaystyle f'(x)=\lim_{n\to\infty}f'_n(x)$.
\end{ti}

\textbf{证明}:由$Newton-Leibniz$公式,$f_n(x)=f_n(a)+\displaystyle\int_a^x f'_n(t)\,\mathrm{d}t$.设$\{f'_n\}$一致收敛于$g$,由逐项积分,$\displaystyle\int_a^x f'_n(t)\,\mathrm{d}t\to\int_a^x g(t)\,\mathrm{d}t$;又$f_n(a)\to f(a)$,故在上式中取极限得$f(x)=f(a)+\displaystyle\int_a^x g(t)\,\mathrm{d}t$.由变上限积分定理,$f$可导且$f'(x)=g(x)=\displaystyle\lim_{n\to\infty}f'_n(x)$.

\begin{ti}
\textbf{定理6.22(Cauchy-Hadamard)}  对幂级数$\displaystyle\sum_{n=0}^{\infty}a_n x^n$,记$\rho=\displaystyle\lim_{n\to\infty}\sqrt[n]{\left\lvert a_n\right\rvert}$,则\\
(1)$\rho=0$时,幂级数在$(-\infty,\infty)$中绝对收敛;\\
(2)$\rho=+\infty$时,幂级数仅在$x=0$处收敛;\\
(3)$0<\rho<\infty$时,幂级数在$\left\lvert x\right\rvert<\dfrac{1}{\rho}$中绝对收敛,在$\left\lvert x\right\rvert>\dfrac{1}{\rho}$之外发散.此时,称$\dfrac{1}{\rho}$为\textbf{收敛半径}.
\end{ti}

\textbf{证明}:注意到$\displaystyle\lim_{n\to\infty}\sqrt[n]{|a_n x^n|}=\rho|x|$,由数项级数的根值判别法即可得到定理结论.

\begin{ti}
\textbf{定理6.23(幂级数的性质)}  设幂级数$\displaystyle\sum_{n=0}^{\infty}a_n x^n$的收敛半径为$R>0$,则其和函数在$(-R,R)$内连续,
且可在$(-R,R)$内逐项求导与逐项积分任意多次,所得幂级数具有相同的收敛半径.
\end{ti}

\textbf{证明}:对任意闭区间$[-r,r]\subset(-R,R)$,取$\rho$使$r<\rho<R$,则当$|x|\leq r$时$|a_nx^n|\leq|a_n|\rho^n$,而$\displaystyle\sum|a_n|\rho^n$收敛,由$Weierstrass$判别法知幂级数在$[-r,r]$上一致收敛,从而和函数在$(-R,R)$内连续.逐项求导后,新幂级数$\displaystyle\sum na_nx^{n-1}$的系数满足$\displaystyle\lim_{n\to\infty}\sqrt[n]{n|a_n|}=\lim_{n\to\infty}\sqrt[n]{|a_n|}$(因为$\lim\sqrt[n]{n}=1$),故收敛半径不变;逐项积分同理.反复应用即得结论.

\begin{ti}
\textbf{定理6.24(Abel连续性定理)}  设幂级数$\displaystyle\sum_{n=0}^{\infty}a_n x^n$的收敛半径为$R\,(0<R<\infty)$.如果$\displaystyle\sum_{n=0}^{\infty}a_n R^n$收敛,则
$$\lim_{x\to R^-}\sum_{n=0}^{\infty}a_n x^n=\sum_{n=0}^{\infty}a_n R^n;$$
如果$\displaystyle\sum_{n=0}^{\infty}a_n(-R)^n$收敛,则
$$\lim_{x\to -R^+}\sum_{n=0}^{\infty}a_n x^n=\sum_{n=0}^{\infty}a_n(-R)^n.$$
\end{ti}

\textbf{证明}:设$\displaystyle\sum a_nR^n$收敛.当$x\in[0,R]$时,$\displaystyle\sum a_nx^n=\sum a_nR^n\left(\frac{x}{R}\right)^n$,其中$\left(\dfrac{x}{R}\right)^n$关于$n$单调且$\leq 1$,由数项级数的$Abel$判别法,$\displaystyle\sum a_nx^n$在$[0,R]$上一致收敛,故$\displaystyle\lim_{x\to R^-}\sum a_nx^n=\sum a_nR^n$.对$x\in[-R,0]$的情形同理,即得另一等式.

\section{Fourier分析}
\begin{ti}
\textbf{定义7.1(Fourier系数与Fourier级数)}  设$f$在$[-\pi,\pi]$中Riemann可积或广义绝对可积.令
$$a_0=\frac{1}{\pi}\int_{-\pi}^{\pi}f(x)\,\mathrm{d}x,\qquad a_n=\frac{1}{\pi}\int_{-\pi}^{\pi}f(x)\cos(nx)\,\mathrm{d}x,\qquad b_n=\frac{1}{\pi}\int_{-\pi}^{\pi}f(x)\sin(nx)\,\mathrm{d}x,\quad n=1,2,\ldots,$$
$a_0,a_n,b_n$称为$f$的\textbf{Fourier系数},形式和
$$\frac{a_0}{2}+\sum_{n=1}^{\infty}\bigl(a_n\cos(nx)+b_n\sin(nx)\bigr)$$
称为$f$的\textbf{Fourier级数}或Fourier展开,记为$\displaystyle f(x)\sim\frac{a_0}{2}+\sum_{n=1}^{\infty}\bigl(a_n\cos(nx)+b_n\sin(nx)\bigr)$.
\end{ti}

\begin{ti}
\textbf{定理7.2(正交性)}  三角函数系$\{1,\cos x,\sin x,\cos 2x,\sin 2x,\ldots\}$在$[-\pi,\pi]$上两两正交,即对任意的$m,n$,有
$$\int_{-\pi}^{\pi}\cos(mx)\sin(nx)\,\mathrm{d}x=0$$
且当$m\neq n$时,有
$$\int_{-\pi}^{\pi}\cos(mx)\cos(nx)\,\mathrm{d}x=\int_{-\pi}^{\pi}\sin(mx)\sin(nx)\,\mathrm{d}x=0$$.
\end{ti}

\textbf{证明}:由直接计算,
$$\int_{-\pi}^{\pi}\cos(mx)\sin(nx)\,\mathrm{d}x=\frac{1}{2}\int_{-\pi}^{\pi}\left[\sin((m+n)x)+\sin((n-m)x)\right]\mathrm{d}x=0;$$
当$m\neq n$时,
$$\int_{-\pi}^{\pi}\cos(mx)\cos(nx)\,\mathrm{d}x=\frac{1}{2}\int_{-\pi}^{\pi}\left[\cos((m+n)x)+\cos((m-n)x)\right]\mathrm{d}x=0,$$
$\sin(mx)\sin(nx)$的积分同理为零.

\begin{ti}
\textbf{定理7.3(Bessel不等式)}  设函数$f$在$[-\pi,\pi]$上平方可积,则
$$\frac{a_0^2}{2}+\sum_{n=1}^{\infty}\bigl(a_n^2+b_n^2\bigr)\leq\frac{1}{\pi}\int_{-\pi}^{\pi}f(x)^2\,\mathrm{d}x$$.
\end{ti}

\textbf{证明}:设$S_n(x)=\dfrac{a_0}{2}+\displaystyle\sum_{k=1}^{n}(a_k\cos kx+b_k\sin kx)$为Fourier级数的部分和.由三角函数系的正交性,
$$\frac{1}{\pi}\int_{-\pi}^{\pi}\left[f(x)-S_n(x)\right]^2\,\mathrm{d}x=\frac{1}{\pi}\int_{-\pi}^{\pi}f(x)^2\,\mathrm{d}x-\left[\frac{a_0^2}{2}+\sum_{k=1}^{n}(a_k^2+b_k^2)\right]\geq 0,$$
令$n\to\infty$即得$Bessel$不等式.

\begin{ti}
\textbf{定理7.4(Riemann-Lebesgue)}  设$f$在$[a,b]$中可积或广义绝对可积,则
$$\lim_{\lambda\to\infty}\int_a^b f(x)\cos(\lambda x)\,\mathrm{d}x=\lim_{\lambda\to\infty}\int_a^b f(x)\sin(\lambda x)\,\mathrm{d}x=0.$$
\end{ti}

\textbf{证明}:先设$f$为阶梯函数或分段线性函数,则通过直接计算(分部积分)可知$\displaystyle\int_a^b f(x)\cos(\lambda x)\,\mathrm{d}x=O\left(\frac{1}{\lambda}\right)$,$\displaystyle\int_a^b f(x)\sin(\lambda x)\,\mathrm{d}x=O\left(\frac{1}{\lambda}\right)$,故极限为$0$.对一般的可积函数$f$,用阶梯函数(或分段线性函数)逼近$f$,使$\sup|f-\varphi|$任意小,由三角不等式即可完成证明.对广义绝对可积函数,用可积函数逼近即可.

\begin{ti}
\textbf{定理7.5(Dirichlet收敛定理)}  设函数$f$以$2\pi$为周期且分段光滑,则其Fourier级数在每一点$x$处收敛到
$$\frac{f(x+0)+f(x-0)}{2}$$
特别地,在$f$的连续点处,Fourier级数收敛到$f(x)$.
\end{ti}

\textbf{证明}:由$Dirichlet$核的计算,Fourier级数的部分和可写为
$$S_n(x)=\frac{1}{\pi}\int_0^{\pi}\frac{f(x+u)+f(x-u)}{2}\cdot\frac{\sin\left(n+\frac{1}{2}\right)u}{2\sin\frac{u}{2}}\,\mathrm{d}u.$$
由$Riemann$局部化原理,部分和的收敛性只与$f$在$x$附近的性态有关.考虑差
$$\frac{f(x+u)+f(x-u)}{2}-\frac{f(x+0)+f(x-0)}{2},$$
当$f$分段光滑时,$u\mapsto\dfrac{f(x+u)+f(x-u)}{2}-\dfrac{f(x+0)+f(x-0)}{2}$除以$\sin\dfrac{u}{2}$后在$u=0$附近为可积函数,利用$Riemann-Lebesgue$引理可得$S_n(x)\to\dfrac{f(x+0)+f(x-0)}{2}$.特别地,在$f$的连续点处收敛到$f(x)$.

\begin{ti}
\textbf{定理7.6(一致收敛)}  设函数$f$以$2\pi$为周期且连续,并具有分段连续的一阶导数,则其Fourier级数在$\mathbb{R}$上一致收敛到$f$.
\end{ti}

\textbf{证明}:设$f$的Fourier系数为$a_n,b_n$,$f'$的Fourier系数为$a'_n,b'_n$.因$f$连续,分部积分无边界项,故$a'_n=nb_n$,$b'_n=-na_n$.由$Riemann-Lebesgue$引理,$a'_n,b'_n=o(1)$,故$|a_n|+|b_n|=o\left(\dfrac{1}{n}\right)$,从而$\displaystyle\sum(|a_n|+|b_n|)$收敛.由$Weierstrass$判别法,$f$的Fourier级数在$\mathbb{R}$上一致收敛;由$Dirichlet$收敛定理其和函数为$f$,故其Fourier级数一致收敛到$f$.

\begin{ti}
\textbf{定理7.7(Parseval恒等式)}  设$f\in R[-\pi,\pi]$,其Fourier展开为$f(x)\sim\dfrac{a_0}{2}+\displaystyle\sum_{n=1}^{\infty}\bigl(a_n\cos(nx)+b_n\sin(nx)\bigr)$,则
$$\frac{1}{\pi}\int_{-\pi}^{\pi}f(x)^2\,\mathrm{d}x=\frac{a_0^2}{2}+\sum_{n=1}^{\infty}\bigl(a_n^2+b_n^2\bigr)$$.
\end{ti}

\textbf{证明}:由$Bessel$不等式的推导,任给$n\geq 1$,
$$\frac{1}{\pi}\int_{-\pi}^{\pi}\left[f(x)-S_n(x)\right]^2\,\mathrm{d}x=\frac{1}{\pi}\int_{-\pi}^{\pi}f(x)^2\,\mathrm{d}x-\left[\frac{a_0^2}{2}+\sum_{k=1}^{n}(a_k^2+b_k^2)\right].$$
利用分段线性函数逼近$f$(必要时先调整端点的函数值使$f(-\pi)=f(\pi)$),可以证明$\dfrac{1}{\pi}\displaystyle\int_{-\pi}^{\pi}[f(x)-S_n(x)]^2\,\mathrm{d}x\to 0$,由此即得$Parseval$恒等式.

\begin{ti}
\textbf{定理7.8(Fourier系数的唯一性)}  设函数$f,g$在$[-\pi,\pi]$上连续且具有相同的Fourier系数,则$f\equiv g$.
\end{ti}

\textbf{证明}:设$h=f-g$,则$h\in C^0[-\pi,\pi]$且其Fourier系数全为零.由$Parseval$恒等式,$\dfrac{1}{\pi}\displaystyle\int_{-\pi}^{\pi}h(x)^2\,\mathrm{d}x=0$.因$h^2$为连续非负函数,故$h\equiv 0$,即$f\equiv g$.
\BookEnd
